% Section LaTeX a inclure dans le memoire avec :
% % Section LaTeX a inclure dans le memoire avec :
% % Section LaTeX a inclure dans le memoire avec :
% % Section LaTeX a inclure dans le memoire avec :
% \input{docs/memoire/devsecops-diagrammes-pas-a-pas}
%
% Packages requis dans le preambule du document principal :
% \usepackage{tikz}
% \usepackage{graphicx}
% \usepackage{array}
% \usetikzlibrary{arrows.meta,positioning,shapes.geometric,fit,backgrounds}

\section{Description complete de la partie DevSecOps}
\label{sec:devsecops-diagrammes-pas-a-pas}

Cette section presente la partie DevSecOps de \textit{SecureRAG Hub} sous la
forme d'un enchainement de pas techniques. L'objectif est de montrer comment le
projet passe d'un code applicatif Laravel versionne a une plateforme
production-like deployee sur Kubernetes, securisee, auditee et documentee par
des preuves.

Le perimetre de cette description est volontairement limite a la chaine
DevSecOps, a Kubernetes, a la supply chain, aux controles de securite, a
l'observabilite, aux preuves runtime et a la documentation. Les composants
metiers hors DevSecOps ne sont pas detailles ici.

\subsection{Creation progressive du socle DevSecOps}
\label{subsec:creation-socle-devsecops}

La creation de la partie DevSecOps n'a pas ete abordee comme un simple ajout
d'outils autour du projet. Elle a ete construite comme un systeme complet de
livraison, de securisation et de preuve. La demarche a consiste a partir du
besoin de soutenance et d'exploitation locale, puis a transformer ce besoin en
une chaine technique reproductible.

Le principe directeur a ete le suivant : chaque automatisation ajoutee au
projet devait produire soit un artefact deployable, soit un rapport de
validation, soit une preuve exploitable dans le memoire. Cette regle a permis
d'eviter une accumulation d'outils non relies entre eux.

\subsubsection{Objectifs initiaux de creation}

Les objectifs initiaux de la creation DevSecOps etaient les suivants :

\begin{itemize}
  \item disposer d'un depot organise et lisible ;
  \item automatiser les tests et les controles de securite ;
  \item construire des images OCI de maniere reproductible ;
  \item deployer localement sur Kubernetes sans infrastructure cloud ;
  \item separer les environnements \texttt{dev}, \texttt{demo} et
  \texttt{production} ;
  \item produire des preuves techniques reutilisables dans la soutenance ;
  \item distinguer honnetement les controles termines des controles dependants
  de l'environnement.
\end{itemize}

\subsubsection{Construction par couches}

La creation s'est faite par couches successives. Chaque couche ajoute un niveau
de maturite supplementaire sans remettre en cause les couches precedentes.

\begin{figure}[htbp]
\centering
\resizebox{\textwidth}{!}{%
\begin{tikzpicture}[node distance=9mm and 8mm]
  \node[srStep] (repo) {Couche 1\\Depot et structure};
  \node[srStep, right=of repo] (ci) {Couche 2\\CI et tests};
  \node[srSec, right=of ci] (security) {Couche 3\\Scans securite};
  \node[srStep, right=of security] (images) {Couche 4\\Images OCI};
  \node[srRuntime, right=of images] (k8s) {Couche 5\\Kubernetes};
  \node[srSec, right=of k8s] (supply) {Couche 6\\Supply chain};
  \node[srProof, right=of supply] (proofs) {Couche 7\\Preuves};

  \draw[srArrow] (repo) -- (ci);
  \draw[srArrow] (ci) -- (security);
  \draw[srArrow] (security) -- (images);
  \draw[srArrow] (images) -- (k8s);
  \draw[srArrow] (k8s) -- (supply);
  \draw[srArrow] (supply) -- (proofs);
\end{tikzpicture}}
\caption{Creation progressive du socle DevSecOps par couches}
\label{fig:creation-devsecops-couches}
\end{figure}

La premiere couche a consiste a stabiliser l'organisation du depot. Les
dossiers \texttt{scripts/}, \texttt{infra/}, \texttt{docs/}, \texttt{security/},
\texttt{platform/} et \texttt{services-laravel/} donnent une responsabilite
claire a chaque partie. Cette separation est importante car elle rend possible
l'automatisation sans melanger code applicatif, infrastructure, runbooks et
preuves.

La deuxieme couche a introduit les scripts de CI. L'objectif etait d'obtenir un
premier signal de qualite avant de construire les images. Les controles de base
sont regroupes dans \texttt{scripts/ci/} et exposes dans le \texttt{Makefile}.
Cette couche verifie notamment les tests, les audits de dependances, les
regles Sonar et les validations statiques.

La troisieme couche a ajoute les controles de securite. Les outils comme
\texttt{Semgrep}, \texttt{Gitleaks} et \texttt{Trivy} permettent de detecter
des erreurs avant la phase de deploiement. Le choix methodologique est de
placer la securite le plus tot possible dans la chaine, afin d'eviter que les
faiblesses ne soient decouvertes uniquement en runtime.

La quatrieme couche a porte sur la construction des images. Les Dockerfiles et
les scripts de build ont ete alignes sur le runtime Laravel officiel. Les
anciens composants legacy sans sources applicatives exploitables ne sont pas
presentes comme build officiel. Ce point est essentiel pour garder une chaine
de preuve honnete.

La cinquieme couche a introduit Kubernetes avec \texttt{kind} et Kustomize. Le
choix de \texttt{kind} permet de simuler localement un environnement
Kubernetes, tandis que Kustomize permet de separer la base commune des overlays
\texttt{dev}, \texttt{demo} et \texttt{production}.

La sixieme couche a ajoute la supply chain : generation de SBOM, scan d'image,
signature Cosign, verification, promotion par digest et deploiement sans
reconstruction. Cette couche transforme le projet d'un simple deploiement local
vers une logique de confiance sur les artefacts.

La septieme couche a formalise la preuve. Les scripts de validation et les
support packs produisent des fichiers sous \texttt{artifacts/}. Ces fichiers
servent a justifier l'etat reel du projet et evitent de declarer un controle
comme termine sans execution correspondante.

\subsubsection{Decisions de conception DevSecOps}

Plusieurs decisions structurantes ont ete prises pendant la creation :

\begin{itemize}
  \item \textbf{Jenkins comme autorite CI/CD} : Jenkins est retenu comme chaine
  officielle pour eviter une double source de verite avec GitHub Actions.
  \item \textbf{Kustomize plutot que manifests dupliques} : la base Kubernetes
  est commune et les differences sont portees par des overlays.
  \item \textbf{Registre local} : \texttt{localhost:5001} permet de demontrer
  la chaine complete sans dependance cloud.
  \item \textbf{Promotion sans rebuild} : une image validee ne doit pas etre
  reconstruite au moment du deploiement.
  \item \textbf{Preuves avant affirmation} : un controle est declare termine
  uniquement si un rapport ou une sortie runtime le prouve.
  \item \textbf{Separation legacy/officiel} : le runtime officiel ne contient
  que les composants Laravel reellement buildables et deployables.
\end{itemize}

\subsubsection{Artefacts crees pendant la construction}

\begin{center}
\renewcommand{\arraystretch}{1.2}
\begin{tabular}{|p{38mm}|p{65mm}|p{42mm}|}
\hline
\textbf{Famille} & \textbf{Role dans la creation DevSecOps} & \textbf{Exemples} \\
\hline
Scripts CI & Executer tests, audits et validations statiques & \texttt{scripts/ci/} \\
\hline
Scripts release & Produire SBOM, signatures, verification et promotion & \texttt{scripts/release/} \\
\hline
Scripts deploy & Creer le cluster, construire les images et deployer & \texttt{scripts/deploy/} \\
\hline
Scripts validate & Collecter les preuves runtime et les rapports & \texttt{scripts/validate/} \\
\hline
Infrastructure & Decrire Jenkins, kind et Kubernetes & \texttt{infra/} \\
\hline
Documentation & Expliquer l'exploitation et les limites reelles & \texttt{docs/runbooks/} \\
\hline
Preuves & Archiver les resultats et statuts reels & \texttt{artifacts/} \\
\hline
\end{tabular}
\end{center}

\subsubsection{Enchainement de creation operationnelle}

L'enchainement operationnel de creation peut etre resume ainsi :

\begin{enumerate}
  \item creation de l'arborescence technique du depot ;
  \item ajout du \texttt{Makefile} comme interface de commandes ;
  \item ajout des scripts CI pour tester et verifier ;
  \item ajout des Dockerfiles et scripts de build ;
  \item ajout du cluster local \texttt{kind} et du registre local ;
  \item creation des manifests Kubernetes de base ;
  \item creation des overlays \texttt{dev}, \texttt{demo} et
  \texttt{production} ;
  \item ajout des NetworkPolicies, PDB, quotas, limites et probes ;
  \item ajout de la supply chain avec SBOM, signatures et promotion ;
  \item ajout des scripts de validation et de support pack ;
  \item alignement de la documentation avec l'etat reel.
\end{enumerate}

Cette creation progressive donne au projet une trajectoire claire : partir d'un
demonstrateur local stable, puis evoluer vers une plateforme production-like
avec controles de securite, haute disponibilite statique, preuves runtime et
release evidence.

\subsection{Legende des diagrammes}

\tikzset{
  srStep/.style={
    rectangle,
    rounded corners=2pt,
    draw=black,
    thick,
    align=center,
    minimum height=8mm,
    text width=31mm,
    fill=gray!10
  },
  srGate/.style={
    diamond,
    draw=black,
    thick,
    align=center,
    aspect=2,
    text width=24mm,
    fill=yellow!15
  },
  srSec/.style={
    rectangle,
    rounded corners=2pt,
    draw=black,
    thick,
    align=center,
    minimum height=8mm,
    text width=32mm,
    fill=red!10
  },
  srProof/.style={
    rectangle,
    rounded corners=2pt,
    draw=black,
    thick,
    align=center,
    minimum height=8mm,
    text width=35mm,
    fill=green!10
  },
  srRuntime/.style={
    rectangle,
    rounded corners=2pt,
    draw=black,
    thick,
    align=center,
    minimum height=8mm,
    text width=34mm,
    fill=blue!8
  },
  srArrow/.style={-{Latex[length=2mm]}, thick}
}

Dans les figures suivantes, les blocs gris representent les etapes techniques,
les blocs rouges les controles de securite, les losanges les gates bloquants,
les blocs bleus les elements runtime Kubernetes et les blocs verts les preuves
archivees.

\subsection{Vue globale de la chaine DevSecOps}

La chaine globale suit une logique progressive : le code est valide en CI, les
images sont construites et securisees, la release est promue par digest, le
deploiement Kubernetes est applique sans reconstruction, puis les preuves
runtime et release sont collectees.

\begin{figure}[htbp]
\centering
\resizebox{\textwidth}{!}{%
\begin{tikzpicture}[node distance=12mm and 10mm]
  \node[srStep] (git) {Depot Git\\monorepo};
  \node[srStep, right=of git] (ci) {Jenkins CI\\qualite et securite};
  \node[srSec, right=of ci] (scan) {SAST, secrets,\\dependances};
  \node[srStep, right=of scan] (build) {Build images\\OCI};
  \node[srSec, right=of build] (supply) {SBOM, signature,\\verification};
  \node[srGate, right=of supply] (gate) {Release\\gate};
  \node[srRuntime, right=of gate] (k8s) {Kubernetes\\production};
  \node[srProof, right=of k8s] (proof) {Preuves\\runtime et release};

  \draw[srArrow] (git) -- (ci);
  \draw[srArrow] (ci) -- (scan);
  \draw[srArrow] (scan) -- (build);
  \draw[srArrow] (build) -- (supply);
  \draw[srArrow] (supply) -- (gate);
  \draw[srArrow] (gate) -- (k8s);
  \draw[srArrow] (k8s) -- (proof);
\end{tikzpicture}}
\caption{Vue globale de la chaine DevSecOps de SecureRAG Hub}
\label{fig:devsecops-global}
\end{figure}

\subsection{Pas 1 : gouvernance du depot et source de verite}

Le premier pas DevSecOps consiste a rendre le depot lisible et gouvernable. Les
sources applicatives, les scripts, l'infrastructure, les runbooks et les
preuves sont separes dans des dossiers dedies. Jenkins est conserve comme
autorite CI/CD officielle, tandis que GitHub Actions reste un heritage
technique non prioritaire.

\begin{figure}[htbp]
\centering
\resizebox{0.95\textwidth}{!}{%
\begin{tikzpicture}[node distance=10mm and 12mm]
  \node[srStep] (repo) {Monorepo\\SecureRAG Hub};
  \node[srStep, below left=of repo] (app) {platform/\\services-laravel/};
  \node[srStep, below=of repo] (infra) {infra/\\k8s, kind, jenkins};
  \node[srStep, below right=of repo] (scripts) {scripts/\\ci, release, deploy};
  \node[srStep, right=of scripts] (docs) {docs/\\runbooks, memoire};
  \node[srProof, right=of repo] (truth) {Source de verite\\Jenkins + docs};

  \draw[srArrow] (repo) -- (app);
  \draw[srArrow] (repo) -- (infra);
  \draw[srArrow] (repo) -- (scripts);
  \draw[srArrow] (repo) -- (docs);
  \draw[srArrow] (repo) -- (truth);
\end{tikzpicture}}
\caption{Organisation du depot et source de verite DevSecOps}
\label{fig:devsecops-repo}
\end{figure}

\subsection{Pas 2 : integration continue Jenkins}

La CI verifie que le code peut etre livre sans introduire de regression
evidente. Elle execute les tests, les validations statiques, les audits de
dependances et les controles de securite. Les resultats sont archives sous
forme de rapports exploitables.

\begin{figure}[htbp]
\centering
\resizebox{\textwidth}{!}{%
\begin{tikzpicture}[node distance=11mm and 9mm]
  \node[srStep] (push) {Push ou\\merge request};
  \node[srStep, right=of push] (checkout) {Checkout\\Jenkins};
  \node[srStep, right=of checkout] (lint) {Lint et rendu\\Kustomize};
  \node[srStep, right=of lint] (tests) {Tests Laravel\\et coverage};
  \node[srSec, right=of tests] (security) {Semgrep\\Gitleaks\\Trivy FS};
  \node[srGate, right=of security] (gate) {Quality\\Gate CI};
  \node[srProof, right=of gate] (reports) {Rapports CI\\archives};

  \draw[srArrow] (push) -- (checkout);
  \draw[srArrow] (checkout) -- (lint);
  \draw[srArrow] (lint) -- (tests);
  \draw[srArrow] (tests) -- (security);
  \draw[srArrow] (security) -- (gate);
  \draw[srArrow] (gate) -- (reports);
\end{tikzpicture}}
\caption{Pipeline Jenkins CI}
\label{fig:devsecops-ci}
\end{figure}

\subsection{Pas 3 : controles de securite applicative et dependances}

Les controles applicatifs sont positionnes tot dans la chaine. Ils incluent la
recherche de secrets, les scans de dependances Composer et npm lorsque les
lockfiles sont presents, les regles SAST et les validations Sonar lorsque
l'environnement fournit les variables necessaires.

\begin{figure}[htbp]
\centering
\resizebox{0.98\textwidth}{!}{%
\begin{tikzpicture}[node distance=10mm and 10mm]
  \node[srStep] (code) {Code Laravel\\et scripts};
  \node[srSec, below left=of code] (deps) {Composer audit\\npm audit};
  \node[srSec, below=of code] (sast) {Semgrep\\Sonar optionnel};
  \node[srSec, below right=of code] (secrets) {Gitleaks\\secrets};
  \node[srGate, right=of code] (block) {Gate\\bloquant};
  \node[srProof, right=of block] (evidence) {security/reports\\artifacts/ci};

  \draw[srArrow] (code) -- (deps);
  \draw[srArrow] (code) -- (sast);
  \draw[srArrow] (code) -- (secrets);
  \draw[srArrow] (deps) -- (block);
  \draw[srArrow] (sast) -- (block);
  \draw[srArrow] (secrets) -- (block);
  \draw[srArrow] (block) -- (evidence);
\end{tikzpicture}}
\caption{Controles de securite applicative et dependances}
\label{fig:devsecops-security-checks}
\end{figure}

\subsection{Pas 4 : construction des images OCI}

Le build produit les images officielles des workloads Laravel. Le runtime
officiel est separe des dossiers legacy afin de ne pas construire ni deployer
des composants sans source applicative exploitable. Les images sont poussees
dans un registre local pour etre consommees par le cluster kind.

\begin{figure}[htbp]
\centering
\resizebox{\textwidth}{!}{%
\begin{tikzpicture}[node distance=11mm and 9mm]
  \node[srStep] (src) {Sources Laravel\\officielles};
  \node[srStep, right=of src] (dockerfile) {Dockerfiles\\production};
  \node[srStep, right=of dockerfile] (build) {docker build};
  \node[srStep, right=of build] (tag) {Tags\\dev/demo/production};
  \node[srStep, right=of tag] (registry) {Registre local\\localhost:5001};
  \node[srRuntime, right=of registry] (kind) {Cluster kind\\pull images};

  \draw[srArrow] (src) -- (dockerfile);
  \draw[srArrow] (dockerfile) -- (build);
  \draw[srArrow] (build) -- (tag);
  \draw[srArrow] (tag) -- (registry);
  \draw[srArrow] (registry) -- (kind);
\end{tikzpicture}}
\caption{Construction et publication des images OCI}
\label{fig:devsecops-build-images}
\end{figure}

\subsection{Pas 5 : supply chain securisee}

La supply chain vise a prouver qu'une image livree correspond bien a l'artefact
construit, scanne, documente par SBOM, signe et verifie. Le projet suit une
logique \textit{digest-first} pour eviter de reconstruire une image au moment du
deploiement.

\begin{figure}[htbp]
\centering
\resizebox{\textwidth}{!}{%
\begin{tikzpicture}[node distance=10mm and 8mm]
  \node[srStep] (image) {Image candidate};
  \node[srSec, right=of image] (trivy) {Trivy\\image scan};
  \node[srSec, right=of trivy] (sbom) {Syft\\SBOM};
  \node[srSec, right=of sbom] (sign) {Cosign\\sign};
  \node[srSec, right=of sign] (verify) {Cosign\\verify};
  \node[srGate, right=of verify] (gate) {Evidence\\gate};
  \node[srStep, right=of gate] (promote) {Promotion\\par digest};
  \node[srProof, right=of promote] (release) {Attestation\\release};

  \draw[srArrow] (image) -- (trivy);
  \draw[srArrow] (trivy) -- (sbom);
  \draw[srArrow] (sbom) -- (sign);
  \draw[srArrow] (sign) -- (verify);
  \draw[srArrow] (verify) -- (gate);
  \draw[srArrow] (gate) -- (promote);
  \draw[srArrow] (promote) -- (release);
\end{tikzpicture}}
\caption{Supply chain securisee : scan, SBOM, signature, verification et promotion}
\label{fig:devsecops-supply-chain}
\end{figure}

\subsection{Pas 6 : deploiement CD sans reconstruction}

Le CD ne doit pas reconstruire les images. Il verifie les artefacts promus,
applique l'overlay Kustomize choisi et collecte les preuves post-deploiement.
Cette separation limite le risque de divergence entre l'image testee et l'image
deployee.

\begin{figure}[htbp]
\centering
\resizebox{\textwidth}{!}{%
\begin{tikzpicture}[node distance=10mm and 9mm]
  \node[srProof] (digest) {Digest promu\\promotion-digests.txt};
  \node[srGate, right=of digest] (verify) {Verification\\pre-deploy};
  \node[srStep, right=of verify] (kustomize) {Kustomize\\overlay};
  \node[srRuntime, right=of kustomize] (apply) {kubectl apply\\cluster};
  \node[srStep, right=of apply] (smoke) {Smoke tests\\security smoke};
  \node[srProof, right=of smoke] (runtime) {Runtime\\evidence};

  \draw[srArrow] (digest) -- (verify);
  \draw[srArrow] (verify) -- (kustomize);
  \draw[srArrow] (kustomize) -- (apply);
  \draw[srArrow] (apply) -- (smoke);
  \draw[srArrow] (smoke) -- (runtime);
\end{tikzpicture}}
\caption{Deploiement CD sans reconstruction d'image}
\label{fig:devsecops-cd-no-rebuild}
\end{figure}

\subsection{Pas 7 : architecture Kubernetes production-like}

L'overlay \texttt{production} conserve le perimetre officiel Laravel et ajoute
les controles attendus pour une plateforme production-like : plusieurs replicas,
rolling update controle, PDB, anti-affinity, topology spread constraints, HPA,
NetworkPolicies et exposition officielle du portail par un Service Kubernetes.

\begin{figure}[htbp]
\centering
\resizebox{\textwidth}{!}{%
\begin{tikzpicture}[node distance=10mm and 10mm]
  \node[srRuntime] (ns) {Namespace\\securerag-hub};
  \node[srRuntime, below left=of ns] (portal) {portal-web\\replicas 3};
  \node[srRuntime, below=of ns] (services) {Services Laravel\\replicas 2};
  \node[srRuntime, below right=of ns] (network) {NetworkPolicies\\default deny};
  \node[srSec, right=of ns] (hardening) {PSA restricted\\securityContext};
  \node[srRuntime, right=of hardening] (ha) {PDB, HPA\\spread};
  \node[srProof, right=of ha] (proof) {production\\runtime evidence};

  \draw[srArrow] (ns) -- (portal);
  \draw[srArrow] (ns) -- (services);
  \draw[srArrow] (ns) -- (network);
  \draw[srArrow] (ns) -- (hardening);
  \draw[srArrow] (hardening) -- (ha);
  \draw[srArrow] (ha) -- (proof);
\end{tikzpicture}}
\caption{Overlay Kubernetes production-like}
\label{fig:devsecops-k8s-production}
\end{figure}

\subsection{Pas 8 : metrics-server et HPA}

Le HPA n'est credible que si \texttt{metrics-server} est disponible et si les
valeurs CPU/memoire ne restent pas a \texttt{\string<unknown\string>}. Le projet distingue
donc les HPA rendus statiquement et les HPA prouves runtime.

\begin{figure}[htbp]
\centering
\resizebox{0.98\textwidth}{!}{%
\begin{tikzpicture}[node distance=10mm and 10mm]
  \node[srRuntime] (cluster) {Cluster\\actif};
  \node[srRuntime, right=of cluster] (metrics) {metrics-server\\Available};
  \node[srStep, right=of metrics] (top) {kubectl top\\nodes/pods};
  \node[srRuntime, right=of top] (hpa) {HPA\\targets remplis};
  \node[srGate, right=of hpa] (gate) {Aucun\\\texttt{\string<unknown\string>}};
  \node[srProof, right=of gate] (report) {hpa-runtime\\report};

  \draw[srArrow] (cluster) -- (metrics);
  \draw[srArrow] (metrics) -- (top);
  \draw[srArrow] (top) -- (hpa);
  \draw[srArrow] (hpa) -- (gate);
  \draw[srArrow] (gate) -- (report);
\end{tikzpicture}}
\caption{Validation runtime de metrics-server et des HPA}
\label{fig:devsecops-hpa}
\end{figure}

\subsection{Pas 9 : Kyverno Audit puis preparation Enforce}

Kyverno est utilise pour auditer les exigences de securite Kubernetes :
profil pod restreint, absence de privileges, references d'images, exposition
des Services et verification Cosign. Le mode \texttt{Enforce} n'est prepare
qu'apres preuve de la supply chain.

\begin{figure}[htbp]
\centering
\resizebox{\textwidth}{!}{%
\begin{tikzpicture}[node distance=10mm and 9mm]
  \node[srRuntime] (cluster) {Cluster\\Kubernetes};
  \node[srSec, right=of cluster] (kyverno) {Kyverno\\Audit};
  \node[srStep, right=of kyverno] (policies) {ClusterPolicies\\SecureRAG};
  \node[srProof, right=of policies] (reports) {PolicyReports};
  \node[srGate, right=of reports] (ready) {Supply chain\\prouvee ?};
  \node[srSec, right=of ready] (enforce) {Enforce\\prudent};

  \draw[srArrow] (cluster) -- (kyverno);
  \draw[srArrow] (kyverno) -- (policies);
  \draw[srArrow] (policies) -- (reports);
  \draw[srArrow] (reports) -- (ready);
  \draw[srArrow] (ready) -- node[above]{oui} (enforce);
\end{tikzpicture}}
\caption{Passage progressif de Kyverno Audit vers Enforce}
\label{fig:devsecops-kyverno}
\end{figure}

\subsection{Pas 10 : gestion des secrets}

La gestion des secrets repose sur une separation entre les valeurs exemples,
les secrets locaux generes, les credentials Jenkins et les secrets Kubernetes.
Aucun secret reel ne doit etre versionne. La rotation reste documentee et doit
etre prouvee dans un environnement d'exploitation.

\begin{figure}[htbp]
\centering
\resizebox{0.98\textwidth}{!}{%
\begin{tikzpicture}[node distance=10mm and 10mm]
  \node[srSec] (strategy) {Strategie\\secrets};
  \node[srStep, below left=of strategy] (examples) {.env.example\\sans secret reel};
  \node[srStep, below=of strategy] (local) {Bootstrap\\local};
  \node[srStep, below right=of strategy] (jenkins) {Credentials\\Jenkins};
  \node[srRuntime, right=of local] (k8s) {Secrets\\Kubernetes};
  \node[srProof, right=of k8s] (audit) {Runbook\\rotation};

  \draw[srArrow] (strategy) -- (examples);
  \draw[srArrow] (strategy) -- (local);
  \draw[srArrow] (strategy) -- (jenkins);
  \draw[srArrow] (local) -- (k8s);
  \draw[srArrow] (jenkins) -- (k8s);
  \draw[srArrow] (k8s) -- (audit);
\end{tikzpicture}}
\caption{Gestion des secrets DevSecOps}
\label{fig:devsecops-secrets}
\end{figure}

\subsection{Pas 11 : resilience donnees, backup et restore}

La resilience donnees production-like impose de sortir d'une logique de
stockage local fragile. L'approche cible est d'utiliser une base externe ou
dediee, de produire des backups, de tester le restore et d'archiver la preuve.

\begin{figure}[htbp]
\centering
\resizebox{\textwidth}{!}{%
\begin{tikzpicture}[node distance=10mm and 9mm]
  \node[srRuntime] (app) {Workloads\\Laravel};
  \node[srRuntime, right=of app] (db) {DB externe\\PostgreSQL};
  \node[srStep, right=of db] (backup) {Backup\\planifie};
  \node[srGate, right=of backup] (restore) {Restore\\teste};
  \node[srProof, right=of restore] (evidence) {Preuve\\data resilience};

  \draw[srArrow] (app) -- (db);
  \draw[srArrow] (db) -- (backup);
  \draw[srArrow] (backup) -- (restore);
  \draw[srArrow] (restore) -- (evidence);
\end{tikzpicture}}
\caption{Resilience donnees : DB externe, backup et restore}
\label{fig:devsecops-data-resilience}
\end{figure}

\subsection{Pas 12 : observabilite runtime}

L'observabilite minimale exploitable repose sur les objets Kubernetes, les
events, les logs applicatifs, les HPA, les metriques \texttt{kubectl top}, les
rapports Kyverno et les snapshots collectes dans les artefacts. Une stack
Prometheus/Grafana/Loki reste une evolution possible, mais n'est pas requise
pour ne pas surcharger la demonstration.

\begin{figure}[htbp]
\centering
\resizebox{\textwidth}{!}{%
\begin{tikzpicture}[node distance=10mm and 8mm]
  \node[srRuntime] (cluster) {Cluster\\runtime};
  \node[srStep, right=of cluster] (objects) {deploy, pods\\svc, pdb, hpa};
  \node[srStep, right=of objects] (events) {Events\\logs};
  \node[srStep, right=of events] (metrics) {kubectl top\\metrics};
  \node[srStep, right=of metrics] (kyverno) {Kyverno\\reports};
  \node[srProof, right=of kyverno] (snapshot) {Observability\\snapshot};

  \draw[srArrow] (cluster) -- (objects);
  \draw[srArrow] (objects) -- (events);
  \draw[srArrow] (events) -- (metrics);
  \draw[srArrow] (metrics) -- (kyverno);
  \draw[srArrow] (kyverno) -- (snapshot);
\end{tikzpicture}}
\caption{Observabilite runtime exploitable}
\label{fig:devsecops-observability}
\end{figure}

\subsection{Pas 13 : support pack et preuves finales}

La derniere etape consiste a consolider les preuves produites par la CI, la
release, Kubernetes, Kyverno, metrics-server, les HPA et les runbooks dans un
support pack final. Ce support pack evite de presenter oralement des controles
non prouves.

\begin{figure}[htbp]
\centering
\resizebox{\textwidth}{!}{%
\begin{tikzpicture}[node distance=10mm and 8mm]
  \node[srProof] (ci) {Preuves CI};
  \node[srProof, right=of ci] (release) {Preuves\\release};
  \node[srProof, right=of release] (runtime) {Preuves\\runtime};
  \node[srProof, right=of runtime] (security) {Preuves\\securite};
  \node[srStep, right=of security] (pack) {Support\\pack};
  \node[srGate, right=of pack] (truth) {Tableau\\etat final};

  \draw[srArrow] (ci) -- (release);
  \draw[srArrow] (release) -- (runtime);
  \draw[srArrow] (runtime) -- (security);
  \draw[srArrow] (security) -- (pack);
  \draw[srArrow] (pack) -- (truth);
\end{tikzpicture}}
\caption{Consolidation des preuves dans le support pack}
\label{fig:devsecops-support-pack}
\end{figure}

\subsection{Synthese des pas DevSecOps}

\begin{center}
\renewcommand{\arraystretch}{1.2}
\begin{tabular}{|p{8mm}|p{42mm}|p{55mm}|p{35mm}|}
\hline
\textbf{Pas} & \textbf{Objectif} & \textbf{Preuve attendue} & \textbf{Etat cible} \\
\hline
1 & Gouvernance du depot & Structure repo, Jenkins source de verite & Termine \\
\hline
2 & CI Jenkins & Rapports tests, lint, scans & Termine si artefacts presents \\
\hline
3 & Securite code et dependances & Semgrep, Gitleaks, Composer/npm audit, Sonar si configure & Termine ou dependant environnement \\
\hline
4 & Build images & Images OCI officielles dans registre local & Termine si Docker/registry actifs \\
\hline
5 & Supply chain & Trivy, Syft, Cosign sign/verify, digest & Termine seulement avec preuves release \\
\hline
6 & CD sans rebuild & Deploiement d'image promue par digest & Termine avec preuve no-rebuild \\
\hline
7 & Kubernetes production-like & Overlay production, PDB, HPA, NetworkPolicies & Termine statique, runtime a prouver \\
\hline
8 & metrics-server et HPA & kubectl top, HPA sans \texttt{\string<unknown\string>} & Dependant cluster actif \\
\hline
9 & Kyverno Audit/Enforce & CRDs, policies, PolicyReports & Audit d'abord, Enforce prudent \\
\hline
10 & Secrets & Aucun secret reel dans Git, bootstrap controle & Partiel sans rotation prouvee \\
\hline
11 & Data resilience & Backup, restore, preuve minimale & Pret non execute si DB absente \\
\hline
12 & Observabilite & Snapshot runtime, logs, events, metrics & Dependant cluster actif \\
\hline
13 & Support pack & Archive finale des preuves & Termine si artefacts generes \\
\hline
\end{tabular}
\end{center}

\subsection{Formulation recommandee pour le memoire}

La partie DevSecOps de \textit{SecureRAG Hub} ne se limite pas a une execution
CI/CD. Elle constitue une chaine de confiance complete reliant le depot Git, la
CI Jenkins, les controles de securite, la construction d'images OCI, la supply
chain par SBOM et signature, le deploiement Kubernetes production-like, les
policies Kyverno, l'observabilite runtime et l'archivage des preuves. Les
controles ne sont declares termines que lorsque les artefacts correspondants
sont presents. Les controles dependant de Docker, kind, metrics-server, Kyverno,
Cosign, Syft, Trivy, PostgreSQL ou Jenkins sont explicitement distingues afin
d'eviter toute surdeclaration.

%
% Packages requis dans le preambule du document principal :
% \usepackage{tikz}
% \usepackage{graphicx}
% \usepackage{array}
% \usetikzlibrary{arrows.meta,positioning,shapes.geometric,fit,backgrounds}

\section{Description complete de la partie DevSecOps}
\label{sec:devsecops-diagrammes-pas-a-pas}

Cette section presente la partie DevSecOps de \textit{SecureRAG Hub} sous la
forme d'un enchainement de pas techniques. L'objectif est de montrer comment le
projet passe d'un code applicatif Laravel versionne a une plateforme
production-like deployee sur Kubernetes, securisee, auditee et documentee par
des preuves.

Le perimetre de cette description est volontairement limite a la chaine
DevSecOps, a Kubernetes, a la supply chain, aux controles de securite, a
l'observabilite, aux preuves runtime et a la documentation. Les composants
metiers hors DevSecOps ne sont pas detailles ici.

\subsection{Creation progressive du socle DevSecOps}
\label{subsec:creation-socle-devsecops}

La creation de la partie DevSecOps n'a pas ete abordee comme un simple ajout
d'outils autour du projet. Elle a ete construite comme un systeme complet de
livraison, de securisation et de preuve. La demarche a consiste a partir du
besoin de soutenance et d'exploitation locale, puis a transformer ce besoin en
une chaine technique reproductible.

Le principe directeur a ete le suivant : chaque automatisation ajoutee au
projet devait produire soit un artefact deployable, soit un rapport de
validation, soit une preuve exploitable dans le memoire. Cette regle a permis
d'eviter une accumulation d'outils non relies entre eux.

\subsubsection{Objectifs initiaux de creation}

Les objectifs initiaux de la creation DevSecOps etaient les suivants :

\begin{itemize}
  \item disposer d'un depot organise et lisible ;
  \item automatiser les tests et les controles de securite ;
  \item construire des images OCI de maniere reproductible ;
  \item deployer localement sur Kubernetes sans infrastructure cloud ;
  \item separer les environnements \texttt{dev}, \texttt{demo} et
  \texttt{production} ;
  \item produire des preuves techniques reutilisables dans la soutenance ;
  \item distinguer honnetement les controles termines des controles dependants
  de l'environnement.
\end{itemize}

\subsubsection{Construction par couches}

La creation s'est faite par couches successives. Chaque couche ajoute un niveau
de maturite supplementaire sans remettre en cause les couches precedentes.

\begin{figure}[htbp]
\centering
\resizebox{\textwidth}{!}{%
\begin{tikzpicture}[node distance=9mm and 8mm]
  \node[srStep] (repo) {Couche 1\\Depot et structure};
  \node[srStep, right=of repo] (ci) {Couche 2\\CI et tests};
  \node[srSec, right=of ci] (security) {Couche 3\\Scans securite};
  \node[srStep, right=of security] (images) {Couche 4\\Images OCI};
  \node[srRuntime, right=of images] (k8s) {Couche 5\\Kubernetes};
  \node[srSec, right=of k8s] (supply) {Couche 6\\Supply chain};
  \node[srProof, right=of supply] (proofs) {Couche 7\\Preuves};

  \draw[srArrow] (repo) -- (ci);
  \draw[srArrow] (ci) -- (security);
  \draw[srArrow] (security) -- (images);
  \draw[srArrow] (images) -- (k8s);
  \draw[srArrow] (k8s) -- (supply);
  \draw[srArrow] (supply) -- (proofs);
\end{tikzpicture}}
\caption{Creation progressive du socle DevSecOps par couches}
\label{fig:creation-devsecops-couches}
\end{figure}

La premiere couche a consiste a stabiliser l'organisation du depot. Les
dossiers \texttt{scripts/}, \texttt{infra/}, \texttt{docs/}, \texttt{security/},
\texttt{platform/} et \texttt{services-laravel/} donnent une responsabilite
claire a chaque partie. Cette separation est importante car elle rend possible
l'automatisation sans melanger code applicatif, infrastructure, runbooks et
preuves.

La deuxieme couche a introduit les scripts de CI. L'objectif etait d'obtenir un
premier signal de qualite avant de construire les images. Les controles de base
sont regroupes dans \texttt{scripts/ci/} et exposes dans le \texttt{Makefile}.
Cette couche verifie notamment les tests, les audits de dependances, les
regles Sonar et les validations statiques.

La troisieme couche a ajoute les controles de securite. Les outils comme
\texttt{Semgrep}, \texttt{Gitleaks} et \texttt{Trivy} permettent de detecter
des erreurs avant la phase de deploiement. Le choix methodologique est de
placer la securite le plus tot possible dans la chaine, afin d'eviter que les
faiblesses ne soient decouvertes uniquement en runtime.

La quatrieme couche a porte sur la construction des images. Les Dockerfiles et
les scripts de build ont ete alignes sur le runtime Laravel officiel. Les
anciens composants legacy sans sources applicatives exploitables ne sont pas
presentes comme build officiel. Ce point est essentiel pour garder une chaine
de preuve honnete.

La cinquieme couche a introduit Kubernetes avec \texttt{kind} et Kustomize. Le
choix de \texttt{kind} permet de simuler localement un environnement
Kubernetes, tandis que Kustomize permet de separer la base commune des overlays
\texttt{dev}, \texttt{demo} et \texttt{production}.

La sixieme couche a ajoute la supply chain : generation de SBOM, scan d'image,
signature Cosign, verification, promotion par digest et deploiement sans
reconstruction. Cette couche transforme le projet d'un simple deploiement local
vers une logique de confiance sur les artefacts.

La septieme couche a formalise la preuve. Les scripts de validation et les
support packs produisent des fichiers sous \texttt{artifacts/}. Ces fichiers
servent a justifier l'etat reel du projet et evitent de declarer un controle
comme termine sans execution correspondante.

\subsubsection{Decisions de conception DevSecOps}

Plusieurs decisions structurantes ont ete prises pendant la creation :

\begin{itemize}
  \item \textbf{Jenkins comme autorite CI/CD} : Jenkins est retenu comme chaine
  officielle pour eviter une double source de verite avec GitHub Actions.
  \item \textbf{Kustomize plutot que manifests dupliques} : la base Kubernetes
  est commune et les differences sont portees par des overlays.
  \item \textbf{Registre local} : \texttt{localhost:5001} permet de demontrer
  la chaine complete sans dependance cloud.
  \item \textbf{Promotion sans rebuild} : une image validee ne doit pas etre
  reconstruite au moment du deploiement.
  \item \textbf{Preuves avant affirmation} : un controle est declare termine
  uniquement si un rapport ou une sortie runtime le prouve.
  \item \textbf{Separation legacy/officiel} : le runtime officiel ne contient
  que les composants Laravel reellement buildables et deployables.
\end{itemize}

\subsubsection{Artefacts crees pendant la construction}

\begin{center}
\renewcommand{\arraystretch}{1.2}
\begin{tabular}{|p{38mm}|p{65mm}|p{42mm}|}
\hline
\textbf{Famille} & \textbf{Role dans la creation DevSecOps} & \textbf{Exemples} \\
\hline
Scripts CI & Executer tests, audits et validations statiques & \texttt{scripts/ci/} \\
\hline
Scripts release & Produire SBOM, signatures, verification et promotion & \texttt{scripts/release/} \\
\hline
Scripts deploy & Creer le cluster, construire les images et deployer & \texttt{scripts/deploy/} \\
\hline
Scripts validate & Collecter les preuves runtime et les rapports & \texttt{scripts/validate/} \\
\hline
Infrastructure & Decrire Jenkins, kind et Kubernetes & \texttt{infra/} \\
\hline
Documentation & Expliquer l'exploitation et les limites reelles & \texttt{docs/runbooks/} \\
\hline
Preuves & Archiver les resultats et statuts reels & \texttt{artifacts/} \\
\hline
\end{tabular}
\end{center}

\subsubsection{Enchainement de creation operationnelle}

L'enchainement operationnel de creation peut etre resume ainsi :

\begin{enumerate}
  \item creation de l'arborescence technique du depot ;
  \item ajout du \texttt{Makefile} comme interface de commandes ;
  \item ajout des scripts CI pour tester et verifier ;
  \item ajout des Dockerfiles et scripts de build ;
  \item ajout du cluster local \texttt{kind} et du registre local ;
  \item creation des manifests Kubernetes de base ;
  \item creation des overlays \texttt{dev}, \texttt{demo} et
  \texttt{production} ;
  \item ajout des NetworkPolicies, PDB, quotas, limites et probes ;
  \item ajout de la supply chain avec SBOM, signatures et promotion ;
  \item ajout des scripts de validation et de support pack ;
  \item alignement de la documentation avec l'etat reel.
\end{enumerate}

Cette creation progressive donne au projet une trajectoire claire : partir d'un
demonstrateur local stable, puis evoluer vers une plateforme production-like
avec controles de securite, haute disponibilite statique, preuves runtime et
release evidence.

\subsection{Legende des diagrammes}

\tikzset{
  srStep/.style={
    rectangle,
    rounded corners=2pt,
    draw=black,
    thick,
    align=center,
    minimum height=8mm,
    text width=31mm,
    fill=gray!10
  },
  srGate/.style={
    diamond,
    draw=black,
    thick,
    align=center,
    aspect=2,
    text width=24mm,
    fill=yellow!15
  },
  srSec/.style={
    rectangle,
    rounded corners=2pt,
    draw=black,
    thick,
    align=center,
    minimum height=8mm,
    text width=32mm,
    fill=red!10
  },
  srProof/.style={
    rectangle,
    rounded corners=2pt,
    draw=black,
    thick,
    align=center,
    minimum height=8mm,
    text width=35mm,
    fill=green!10
  },
  srRuntime/.style={
    rectangle,
    rounded corners=2pt,
    draw=black,
    thick,
    align=center,
    minimum height=8mm,
    text width=34mm,
    fill=blue!8
  },
  srArrow/.style={-{Latex[length=2mm]}, thick}
}

Dans les figures suivantes, les blocs gris representent les etapes techniques,
les blocs rouges les controles de securite, les losanges les gates bloquants,
les blocs bleus les elements runtime Kubernetes et les blocs verts les preuves
archivees.

\subsection{Vue globale de la chaine DevSecOps}

La chaine globale suit une logique progressive : le code est valide en CI, les
images sont construites et securisees, la release est promue par digest, le
deploiement Kubernetes est applique sans reconstruction, puis les preuves
runtime et release sont collectees.

\begin{figure}[htbp]
\centering
\resizebox{\textwidth}{!}{%
\begin{tikzpicture}[node distance=12mm and 10mm]
  \node[srStep] (git) {Depot Git\\monorepo};
  \node[srStep, right=of git] (ci) {Jenkins CI\\qualite et securite};
  \node[srSec, right=of ci] (scan) {SAST, secrets,\\dependances};
  \node[srStep, right=of scan] (build) {Build images\\OCI};
  \node[srSec, right=of build] (supply) {SBOM, signature,\\verification};
  \node[srGate, right=of supply] (gate) {Release\\gate};
  \node[srRuntime, right=of gate] (k8s) {Kubernetes\\production};
  \node[srProof, right=of k8s] (proof) {Preuves\\runtime et release};

  \draw[srArrow] (git) -- (ci);
  \draw[srArrow] (ci) -- (scan);
  \draw[srArrow] (scan) -- (build);
  \draw[srArrow] (build) -- (supply);
  \draw[srArrow] (supply) -- (gate);
  \draw[srArrow] (gate) -- (k8s);
  \draw[srArrow] (k8s) -- (proof);
\end{tikzpicture}}
\caption{Vue globale de la chaine DevSecOps de SecureRAG Hub}
\label{fig:devsecops-global}
\end{figure}

\subsection{Pas 1 : gouvernance du depot et source de verite}

Le premier pas DevSecOps consiste a rendre le depot lisible et gouvernable. Les
sources applicatives, les scripts, l'infrastructure, les runbooks et les
preuves sont separes dans des dossiers dedies. Jenkins est conserve comme
autorite CI/CD officielle, tandis que GitHub Actions reste un heritage
technique non prioritaire.

\begin{figure}[htbp]
\centering
\resizebox{0.95\textwidth}{!}{%
\begin{tikzpicture}[node distance=10mm and 12mm]
  \node[srStep] (repo) {Monorepo\\SecureRAG Hub};
  \node[srStep, below left=of repo] (app) {platform/\\services-laravel/};
  \node[srStep, below=of repo] (infra) {infra/\\k8s, kind, jenkins};
  \node[srStep, below right=of repo] (scripts) {scripts/\\ci, release, deploy};
  \node[srStep, right=of scripts] (docs) {docs/\\runbooks, memoire};
  \node[srProof, right=of repo] (truth) {Source de verite\\Jenkins + docs};

  \draw[srArrow] (repo) -- (app);
  \draw[srArrow] (repo) -- (infra);
  \draw[srArrow] (repo) -- (scripts);
  \draw[srArrow] (repo) -- (docs);
  \draw[srArrow] (repo) -- (truth);
\end{tikzpicture}}
\caption{Organisation du depot et source de verite DevSecOps}
\label{fig:devsecops-repo}
\end{figure}

\subsection{Pas 2 : integration continue Jenkins}

La CI verifie que le code peut etre livre sans introduire de regression
evidente. Elle execute les tests, les validations statiques, les audits de
dependances et les controles de securite. Les resultats sont archives sous
forme de rapports exploitables.

\begin{figure}[htbp]
\centering
\resizebox{\textwidth}{!}{%
\begin{tikzpicture}[node distance=11mm and 9mm]
  \node[srStep] (push) {Push ou\\merge request};
  \node[srStep, right=of push] (checkout) {Checkout\\Jenkins};
  \node[srStep, right=of checkout] (lint) {Lint et rendu\\Kustomize};
  \node[srStep, right=of lint] (tests) {Tests Laravel\\et coverage};
  \node[srSec, right=of tests] (security) {Semgrep\\Gitleaks\\Trivy FS};
  \node[srGate, right=of security] (gate) {Quality\\Gate CI};
  \node[srProof, right=of gate] (reports) {Rapports CI\\archives};

  \draw[srArrow] (push) -- (checkout);
  \draw[srArrow] (checkout) -- (lint);
  \draw[srArrow] (lint) -- (tests);
  \draw[srArrow] (tests) -- (security);
  \draw[srArrow] (security) -- (gate);
  \draw[srArrow] (gate) -- (reports);
\end{tikzpicture}}
\caption{Pipeline Jenkins CI}
\label{fig:devsecops-ci}
\end{figure}

\subsection{Pas 3 : controles de securite applicative et dependances}

Les controles applicatifs sont positionnes tot dans la chaine. Ils incluent la
recherche de secrets, les scans de dependances Composer et npm lorsque les
lockfiles sont presents, les regles SAST et les validations Sonar lorsque
l'environnement fournit les variables necessaires.

\begin{figure}[htbp]
\centering
\resizebox{0.98\textwidth}{!}{%
\begin{tikzpicture}[node distance=10mm and 10mm]
  \node[srStep] (code) {Code Laravel\\et scripts};
  \node[srSec, below left=of code] (deps) {Composer audit\\npm audit};
  \node[srSec, below=of code] (sast) {Semgrep\\Sonar optionnel};
  \node[srSec, below right=of code] (secrets) {Gitleaks\\secrets};
  \node[srGate, right=of code] (block) {Gate\\bloquant};
  \node[srProof, right=of block] (evidence) {security/reports\\artifacts/ci};

  \draw[srArrow] (code) -- (deps);
  \draw[srArrow] (code) -- (sast);
  \draw[srArrow] (code) -- (secrets);
  \draw[srArrow] (deps) -- (block);
  \draw[srArrow] (sast) -- (block);
  \draw[srArrow] (secrets) -- (block);
  \draw[srArrow] (block) -- (evidence);
\end{tikzpicture}}
\caption{Controles de securite applicative et dependances}
\label{fig:devsecops-security-checks}
\end{figure}

\subsection{Pas 4 : construction des images OCI}

Le build produit les images officielles des workloads Laravel. Le runtime
officiel est separe des dossiers legacy afin de ne pas construire ni deployer
des composants sans source applicative exploitable. Les images sont poussees
dans un registre local pour etre consommees par le cluster kind.

\begin{figure}[htbp]
\centering
\resizebox{\textwidth}{!}{%
\begin{tikzpicture}[node distance=11mm and 9mm]
  \node[srStep] (src) {Sources Laravel\\officielles};
  \node[srStep, right=of src] (dockerfile) {Dockerfiles\\production};
  \node[srStep, right=of dockerfile] (build) {docker build};
  \node[srStep, right=of build] (tag) {Tags\\dev/demo/production};
  \node[srStep, right=of tag] (registry) {Registre local\\localhost:5001};
  \node[srRuntime, right=of registry] (kind) {Cluster kind\\pull images};

  \draw[srArrow] (src) -- (dockerfile);
  \draw[srArrow] (dockerfile) -- (build);
  \draw[srArrow] (build) -- (tag);
  \draw[srArrow] (tag) -- (registry);
  \draw[srArrow] (registry) -- (kind);
\end{tikzpicture}}
\caption{Construction et publication des images OCI}
\label{fig:devsecops-build-images}
\end{figure}

\subsection{Pas 5 : supply chain securisee}

La supply chain vise a prouver qu'une image livree correspond bien a l'artefact
construit, scanne, documente par SBOM, signe et verifie. Le projet suit une
logique \textit{digest-first} pour eviter de reconstruire une image au moment du
deploiement.

\begin{figure}[htbp]
\centering
\resizebox{\textwidth}{!}{%
\begin{tikzpicture}[node distance=10mm and 8mm]
  \node[srStep] (image) {Image candidate};
  \node[srSec, right=of image] (trivy) {Trivy\\image scan};
  \node[srSec, right=of trivy] (sbom) {Syft\\SBOM};
  \node[srSec, right=of sbom] (sign) {Cosign\\sign};
  \node[srSec, right=of sign] (verify) {Cosign\\verify};
  \node[srGate, right=of verify] (gate) {Evidence\\gate};
  \node[srStep, right=of gate] (promote) {Promotion\\par digest};
  \node[srProof, right=of promote] (release) {Attestation\\release};

  \draw[srArrow] (image) -- (trivy);
  \draw[srArrow] (trivy) -- (sbom);
  \draw[srArrow] (sbom) -- (sign);
  \draw[srArrow] (sign) -- (verify);
  \draw[srArrow] (verify) -- (gate);
  \draw[srArrow] (gate) -- (promote);
  \draw[srArrow] (promote) -- (release);
\end{tikzpicture}}
\caption{Supply chain securisee : scan, SBOM, signature, verification et promotion}
\label{fig:devsecops-supply-chain}
\end{figure}

\subsection{Pas 6 : deploiement CD sans reconstruction}

Le CD ne doit pas reconstruire les images. Il verifie les artefacts promus,
applique l'overlay Kustomize choisi et collecte les preuves post-deploiement.
Cette separation limite le risque de divergence entre l'image testee et l'image
deployee.

\begin{figure}[htbp]
\centering
\resizebox{\textwidth}{!}{%
\begin{tikzpicture}[node distance=10mm and 9mm]
  \node[srProof] (digest) {Digest promu\\promotion-digests.txt};
  \node[srGate, right=of digest] (verify) {Verification\\pre-deploy};
  \node[srStep, right=of verify] (kustomize) {Kustomize\\overlay};
  \node[srRuntime, right=of kustomize] (apply) {kubectl apply\\cluster};
  \node[srStep, right=of apply] (smoke) {Smoke tests\\security smoke};
  \node[srProof, right=of smoke] (runtime) {Runtime\\evidence};

  \draw[srArrow] (digest) -- (verify);
  \draw[srArrow] (verify) -- (kustomize);
  \draw[srArrow] (kustomize) -- (apply);
  \draw[srArrow] (apply) -- (smoke);
  \draw[srArrow] (smoke) -- (runtime);
\end{tikzpicture}}
\caption{Deploiement CD sans reconstruction d'image}
\label{fig:devsecops-cd-no-rebuild}
\end{figure}

\subsection{Pas 7 : architecture Kubernetes production-like}

L'overlay \texttt{production} conserve le perimetre officiel Laravel et ajoute
les controles attendus pour une plateforme production-like : plusieurs replicas,
rolling update controle, PDB, anti-affinity, topology spread constraints, HPA,
NetworkPolicies et exposition officielle du portail par un Service Kubernetes.

\begin{figure}[htbp]
\centering
\resizebox{\textwidth}{!}{%
\begin{tikzpicture}[node distance=10mm and 10mm]
  \node[srRuntime] (ns) {Namespace\\securerag-hub};
  \node[srRuntime, below left=of ns] (portal) {portal-web\\replicas 3};
  \node[srRuntime, below=of ns] (services) {Services Laravel\\replicas 2};
  \node[srRuntime, below right=of ns] (network) {NetworkPolicies\\default deny};
  \node[srSec, right=of ns] (hardening) {PSA restricted\\securityContext};
  \node[srRuntime, right=of hardening] (ha) {PDB, HPA\\spread};
  \node[srProof, right=of ha] (proof) {production\\runtime evidence};

  \draw[srArrow] (ns) -- (portal);
  \draw[srArrow] (ns) -- (services);
  \draw[srArrow] (ns) -- (network);
  \draw[srArrow] (ns) -- (hardening);
  \draw[srArrow] (hardening) -- (ha);
  \draw[srArrow] (ha) -- (proof);
\end{tikzpicture}}
\caption{Overlay Kubernetes production-like}
\label{fig:devsecops-k8s-production}
\end{figure}

\subsection{Pas 8 : metrics-server et HPA}

Le HPA n'est credible que si \texttt{metrics-server} est disponible et si les
valeurs CPU/memoire ne restent pas a \texttt{\string<unknown\string>}. Le projet distingue
donc les HPA rendus statiquement et les HPA prouves runtime.

\begin{figure}[htbp]
\centering
\resizebox{0.98\textwidth}{!}{%
\begin{tikzpicture}[node distance=10mm and 10mm]
  \node[srRuntime] (cluster) {Cluster\\actif};
  \node[srRuntime, right=of cluster] (metrics) {metrics-server\\Available};
  \node[srStep, right=of metrics] (top) {kubectl top\\nodes/pods};
  \node[srRuntime, right=of top] (hpa) {HPA\\targets remplis};
  \node[srGate, right=of hpa] (gate) {Aucun\\\texttt{\string<unknown\string>}};
  \node[srProof, right=of gate] (report) {hpa-runtime\\report};

  \draw[srArrow] (cluster) -- (metrics);
  \draw[srArrow] (metrics) -- (top);
  \draw[srArrow] (top) -- (hpa);
  \draw[srArrow] (hpa) -- (gate);
  \draw[srArrow] (gate) -- (report);
\end{tikzpicture}}
\caption{Validation runtime de metrics-server et des HPA}
\label{fig:devsecops-hpa}
\end{figure}

\subsection{Pas 9 : Kyverno Audit puis preparation Enforce}

Kyverno est utilise pour auditer les exigences de securite Kubernetes :
profil pod restreint, absence de privileges, references d'images, exposition
des Services et verification Cosign. Le mode \texttt{Enforce} n'est prepare
qu'apres preuve de la supply chain.

\begin{figure}[htbp]
\centering
\resizebox{\textwidth}{!}{%
\begin{tikzpicture}[node distance=10mm and 9mm]
  \node[srRuntime] (cluster) {Cluster\\Kubernetes};
  \node[srSec, right=of cluster] (kyverno) {Kyverno\\Audit};
  \node[srStep, right=of kyverno] (policies) {ClusterPolicies\\SecureRAG};
  \node[srProof, right=of policies] (reports) {PolicyReports};
  \node[srGate, right=of reports] (ready) {Supply chain\\prouvee ?};
  \node[srSec, right=of ready] (enforce) {Enforce\\prudent};

  \draw[srArrow] (cluster) -- (kyverno);
  \draw[srArrow] (kyverno) -- (policies);
  \draw[srArrow] (policies) -- (reports);
  \draw[srArrow] (reports) -- (ready);
  \draw[srArrow] (ready) -- node[above]{oui} (enforce);
\end{tikzpicture}}
\caption{Passage progressif de Kyverno Audit vers Enforce}
\label{fig:devsecops-kyverno}
\end{figure}

\subsection{Pas 10 : gestion des secrets}

La gestion des secrets repose sur une separation entre les valeurs exemples,
les secrets locaux generes, les credentials Jenkins et les secrets Kubernetes.
Aucun secret reel ne doit etre versionne. La rotation reste documentee et doit
etre prouvee dans un environnement d'exploitation.

\begin{figure}[htbp]
\centering
\resizebox{0.98\textwidth}{!}{%
\begin{tikzpicture}[node distance=10mm and 10mm]
  \node[srSec] (strategy) {Strategie\\secrets};
  \node[srStep, below left=of strategy] (examples) {.env.example\\sans secret reel};
  \node[srStep, below=of strategy] (local) {Bootstrap\\local};
  \node[srStep, below right=of strategy] (jenkins) {Credentials\\Jenkins};
  \node[srRuntime, right=of local] (k8s) {Secrets\\Kubernetes};
  \node[srProof, right=of k8s] (audit) {Runbook\\rotation};

  \draw[srArrow] (strategy) -- (examples);
  \draw[srArrow] (strategy) -- (local);
  \draw[srArrow] (strategy) -- (jenkins);
  \draw[srArrow] (local) -- (k8s);
  \draw[srArrow] (jenkins) -- (k8s);
  \draw[srArrow] (k8s) -- (audit);
\end{tikzpicture}}
\caption{Gestion des secrets DevSecOps}
\label{fig:devsecops-secrets}
\end{figure}

\subsection{Pas 11 : resilience donnees, backup et restore}

La resilience donnees production-like impose de sortir d'une logique de
stockage local fragile. L'approche cible est d'utiliser une base externe ou
dediee, de produire des backups, de tester le restore et d'archiver la preuve.

\begin{figure}[htbp]
\centering
\resizebox{\textwidth}{!}{%
\begin{tikzpicture}[node distance=10mm and 9mm]
  \node[srRuntime] (app) {Workloads\\Laravel};
  \node[srRuntime, right=of app] (db) {DB externe\\PostgreSQL};
  \node[srStep, right=of db] (backup) {Backup\\planifie};
  \node[srGate, right=of backup] (restore) {Restore\\teste};
  \node[srProof, right=of restore] (evidence) {Preuve\\data resilience};

  \draw[srArrow] (app) -- (db);
  \draw[srArrow] (db) -- (backup);
  \draw[srArrow] (backup) -- (restore);
  \draw[srArrow] (restore) -- (evidence);
\end{tikzpicture}}
\caption{Resilience donnees : DB externe, backup et restore}
\label{fig:devsecops-data-resilience}
\end{figure}

\subsection{Pas 12 : observabilite runtime}

L'observabilite minimale exploitable repose sur les objets Kubernetes, les
events, les logs applicatifs, les HPA, les metriques \texttt{kubectl top}, les
rapports Kyverno et les snapshots collectes dans les artefacts. Une stack
Prometheus/Grafana/Loki reste une evolution possible, mais n'est pas requise
pour ne pas surcharger la demonstration.

\begin{figure}[htbp]
\centering
\resizebox{\textwidth}{!}{%
\begin{tikzpicture}[node distance=10mm and 8mm]
  \node[srRuntime] (cluster) {Cluster\\runtime};
  \node[srStep, right=of cluster] (objects) {deploy, pods\\svc, pdb, hpa};
  \node[srStep, right=of objects] (events) {Events\\logs};
  \node[srStep, right=of events] (metrics) {kubectl top\\metrics};
  \node[srStep, right=of metrics] (kyverno) {Kyverno\\reports};
  \node[srProof, right=of kyverno] (snapshot) {Observability\\snapshot};

  \draw[srArrow] (cluster) -- (objects);
  \draw[srArrow] (objects) -- (events);
  \draw[srArrow] (events) -- (metrics);
  \draw[srArrow] (metrics) -- (kyverno);
  \draw[srArrow] (kyverno) -- (snapshot);
\end{tikzpicture}}
\caption{Observabilite runtime exploitable}
\label{fig:devsecops-observability}
\end{figure}

\subsection{Pas 13 : support pack et preuves finales}

La derniere etape consiste a consolider les preuves produites par la CI, la
release, Kubernetes, Kyverno, metrics-server, les HPA et les runbooks dans un
support pack final. Ce support pack evite de presenter oralement des controles
non prouves.

\begin{figure}[htbp]
\centering
\resizebox{\textwidth}{!}{%
\begin{tikzpicture}[node distance=10mm and 8mm]
  \node[srProof] (ci) {Preuves CI};
  \node[srProof, right=of ci] (release) {Preuves\\release};
  \node[srProof, right=of release] (runtime) {Preuves\\runtime};
  \node[srProof, right=of runtime] (security) {Preuves\\securite};
  \node[srStep, right=of security] (pack) {Support\\pack};
  \node[srGate, right=of pack] (truth) {Tableau\\etat final};

  \draw[srArrow] (ci) -- (release);
  \draw[srArrow] (release) -- (runtime);
  \draw[srArrow] (runtime) -- (security);
  \draw[srArrow] (security) -- (pack);
  \draw[srArrow] (pack) -- (truth);
\end{tikzpicture}}
\caption{Consolidation des preuves dans le support pack}
\label{fig:devsecops-support-pack}
\end{figure}

\subsection{Synthese des pas DevSecOps}

\begin{center}
\renewcommand{\arraystretch}{1.2}
\begin{tabular}{|p{8mm}|p{42mm}|p{55mm}|p{35mm}|}
\hline
\textbf{Pas} & \textbf{Objectif} & \textbf{Preuve attendue} & \textbf{Etat cible} \\
\hline
1 & Gouvernance du depot & Structure repo, Jenkins source de verite & Termine \\
\hline
2 & CI Jenkins & Rapports tests, lint, scans & Termine si artefacts presents \\
\hline
3 & Securite code et dependances & Semgrep, Gitleaks, Composer/npm audit, Sonar si configure & Termine ou dependant environnement \\
\hline
4 & Build images & Images OCI officielles dans registre local & Termine si Docker/registry actifs \\
\hline
5 & Supply chain & Trivy, Syft, Cosign sign/verify, digest & Termine seulement avec preuves release \\
\hline
6 & CD sans rebuild & Deploiement d'image promue par digest & Termine avec preuve no-rebuild \\
\hline
7 & Kubernetes production-like & Overlay production, PDB, HPA, NetworkPolicies & Termine statique, runtime a prouver \\
\hline
8 & metrics-server et HPA & kubectl top, HPA sans \texttt{\string<unknown\string>} & Dependant cluster actif \\
\hline
9 & Kyverno Audit/Enforce & CRDs, policies, PolicyReports & Audit d'abord, Enforce prudent \\
\hline
10 & Secrets & Aucun secret reel dans Git, bootstrap controle & Partiel sans rotation prouvee \\
\hline
11 & Data resilience & Backup, restore, preuve minimale & Pret non execute si DB absente \\
\hline
12 & Observabilite & Snapshot runtime, logs, events, metrics & Dependant cluster actif \\
\hline
13 & Support pack & Archive finale des preuves & Termine si artefacts generes \\
\hline
\end{tabular}
\end{center}

\subsection{Formulation recommandee pour le memoire}

La partie DevSecOps de \textit{SecureRAG Hub} ne se limite pas a une execution
CI/CD. Elle constitue une chaine de confiance complete reliant le depot Git, la
CI Jenkins, les controles de securite, la construction d'images OCI, la supply
chain par SBOM et signature, le deploiement Kubernetes production-like, les
policies Kyverno, l'observabilite runtime et l'archivage des preuves. Les
controles ne sont declares termines que lorsque les artefacts correspondants
sont presents. Les controles dependant de Docker, kind, metrics-server, Kyverno,
Cosign, Syft, Trivy, PostgreSQL ou Jenkins sont explicitement distingues afin
d'eviter toute surdeclaration.

%
% Packages requis dans le preambule du document principal :
% \usepackage{tikz}
% \usepackage{graphicx}
% \usepackage{array}
% \usetikzlibrary{arrows.meta,positioning,shapes.geometric,fit,backgrounds}

\section{Description complete de la partie DevSecOps}
\label{sec:devsecops-diagrammes-pas-a-pas}

Cette section presente la partie DevSecOps de \textit{SecureRAG Hub} sous la
forme d'un enchainement de pas techniques. L'objectif est de montrer comment le
projet passe d'un code applicatif Laravel versionne a une plateforme
production-like deployee sur Kubernetes, securisee, auditee et documentee par
des preuves.

Le perimetre de cette description est volontairement limite a la chaine
DevSecOps, a Kubernetes, a la supply chain, aux controles de securite, a
l'observabilite, aux preuves runtime et a la documentation. Les composants
metiers hors DevSecOps ne sont pas detailles ici.

\subsection{Creation progressive du socle DevSecOps}
\label{subsec:creation-socle-devsecops}

La creation de la partie DevSecOps n'a pas ete abordee comme un simple ajout
d'outils autour du projet. Elle a ete construite comme un systeme complet de
livraison, de securisation et de preuve. La demarche a consiste a partir du
besoin de soutenance et d'exploitation locale, puis a transformer ce besoin en
une chaine technique reproductible.

Le principe directeur a ete le suivant : chaque automatisation ajoutee au
projet devait produire soit un artefact deployable, soit un rapport de
validation, soit une preuve exploitable dans le memoire. Cette regle a permis
d'eviter une accumulation d'outils non relies entre eux.

\subsubsection{Objectifs initiaux de creation}

Les objectifs initiaux de la creation DevSecOps etaient les suivants :

\begin{itemize}
  \item disposer d'un depot organise et lisible ;
  \item automatiser les tests et les controles de securite ;
  \item construire des images OCI de maniere reproductible ;
  \item deployer localement sur Kubernetes sans infrastructure cloud ;
  \item separer les environnements \texttt{dev}, \texttt{demo} et
  \texttt{production} ;
  \item produire des preuves techniques reutilisables dans la soutenance ;
  \item distinguer honnetement les controles termines des controles dependants
  de l'environnement.
\end{itemize}

\subsubsection{Construction par couches}

La creation s'est faite par couches successives. Chaque couche ajoute un niveau
de maturite supplementaire sans remettre en cause les couches precedentes.

\begin{figure}[htbp]
\centering
\resizebox{\textwidth}{!}{%
\begin{tikzpicture}[node distance=9mm and 8mm]
  \node[srStep] (repo) {Couche 1\\Depot et structure};
  \node[srStep, right=of repo] (ci) {Couche 2\\CI et tests};
  \node[srSec, right=of ci] (security) {Couche 3\\Scans securite};
  \node[srStep, right=of security] (images) {Couche 4\\Images OCI};
  \node[srRuntime, right=of images] (k8s) {Couche 5\\Kubernetes};
  \node[srSec, right=of k8s] (supply) {Couche 6\\Supply chain};
  \node[srProof, right=of supply] (proofs) {Couche 7\\Preuves};

  \draw[srArrow] (repo) -- (ci);
  \draw[srArrow] (ci) -- (security);
  \draw[srArrow] (security) -- (images);
  \draw[srArrow] (images) -- (k8s);
  \draw[srArrow] (k8s) -- (supply);
  \draw[srArrow] (supply) -- (proofs);
\end{tikzpicture}}
\caption{Creation progressive du socle DevSecOps par couches}
\label{fig:creation-devsecops-couches}
\end{figure}

La premiere couche a consiste a stabiliser l'organisation du depot. Les
dossiers \texttt{scripts/}, \texttt{infra/}, \texttt{docs/}, \texttt{security/},
\texttt{platform/} et \texttt{services-laravel/} donnent une responsabilite
claire a chaque partie. Cette separation est importante car elle rend possible
l'automatisation sans melanger code applicatif, infrastructure, runbooks et
preuves.

La deuxieme couche a introduit les scripts de CI. L'objectif etait d'obtenir un
premier signal de qualite avant de construire les images. Les controles de base
sont regroupes dans \texttt{scripts/ci/} et exposes dans le \texttt{Makefile}.
Cette couche verifie notamment les tests, les audits de dependances, les
regles Sonar et les validations statiques.

La troisieme couche a ajoute les controles de securite. Les outils comme
\texttt{Semgrep}, \texttt{Gitleaks} et \texttt{Trivy} permettent de detecter
des erreurs avant la phase de deploiement. Le choix methodologique est de
placer la securite le plus tot possible dans la chaine, afin d'eviter que les
faiblesses ne soient decouvertes uniquement en runtime.

La quatrieme couche a porte sur la construction des images. Les Dockerfiles et
les scripts de build ont ete alignes sur le runtime Laravel officiel. Les
anciens composants legacy sans sources applicatives exploitables ne sont pas
presentes comme build officiel. Ce point est essentiel pour garder une chaine
de preuve honnete.

La cinquieme couche a introduit Kubernetes avec \texttt{kind} et Kustomize. Le
choix de \texttt{kind} permet de simuler localement un environnement
Kubernetes, tandis que Kustomize permet de separer la base commune des overlays
\texttt{dev}, \texttt{demo} et \texttt{production}.

La sixieme couche a ajoute la supply chain : generation de SBOM, scan d'image,
signature Cosign, verification, promotion par digest et deploiement sans
reconstruction. Cette couche transforme le projet d'un simple deploiement local
vers une logique de confiance sur les artefacts.

La septieme couche a formalise la preuve. Les scripts de validation et les
support packs produisent des fichiers sous \texttt{artifacts/}. Ces fichiers
servent a justifier l'etat reel du projet et evitent de declarer un controle
comme termine sans execution correspondante.

\subsubsection{Decisions de conception DevSecOps}

Plusieurs decisions structurantes ont ete prises pendant la creation :

\begin{itemize}
  \item \textbf{Jenkins comme autorite CI/CD} : Jenkins est retenu comme chaine
  officielle pour eviter une double source de verite avec GitHub Actions.
  \item \textbf{Kustomize plutot que manifests dupliques} : la base Kubernetes
  est commune et les differences sont portees par des overlays.
  \item \textbf{Registre local} : \texttt{localhost:5001} permet de demontrer
  la chaine complete sans dependance cloud.
  \item \textbf{Promotion sans rebuild} : une image validee ne doit pas etre
  reconstruite au moment du deploiement.
  \item \textbf{Preuves avant affirmation} : un controle est declare termine
  uniquement si un rapport ou une sortie runtime le prouve.
  \item \textbf{Separation legacy/officiel} : le runtime officiel ne contient
  que les composants Laravel reellement buildables et deployables.
\end{itemize}

\subsubsection{Artefacts crees pendant la construction}

\begin{center}
\renewcommand{\arraystretch}{1.2}
\begin{tabular}{|p{38mm}|p{65mm}|p{42mm}|}
\hline
\textbf{Famille} & \textbf{Role dans la creation DevSecOps} & \textbf{Exemples} \\
\hline
Scripts CI & Executer tests, audits et validations statiques & \texttt{scripts/ci/} \\
\hline
Scripts release & Produire SBOM, signatures, verification et promotion & \texttt{scripts/release/} \\
\hline
Scripts deploy & Creer le cluster, construire les images et deployer & \texttt{scripts/deploy/} \\
\hline
Scripts validate & Collecter les preuves runtime et les rapports & \texttt{scripts/validate/} \\
\hline
Infrastructure & Decrire Jenkins, kind et Kubernetes & \texttt{infra/} \\
\hline
Documentation & Expliquer l'exploitation et les limites reelles & \texttt{docs/runbooks/} \\
\hline
Preuves & Archiver les resultats et statuts reels & \texttt{artifacts/} \\
\hline
\end{tabular}
\end{center}

\subsubsection{Enchainement de creation operationnelle}

L'enchainement operationnel de creation peut etre resume ainsi :

\begin{enumerate}
  \item creation de l'arborescence technique du depot ;
  \item ajout du \texttt{Makefile} comme interface de commandes ;
  \item ajout des scripts CI pour tester et verifier ;
  \item ajout des Dockerfiles et scripts de build ;
  \item ajout du cluster local \texttt{kind} et du registre local ;
  \item creation des manifests Kubernetes de base ;
  \item creation des overlays \texttt{dev}, \texttt{demo} et
  \texttt{production} ;
  \item ajout des NetworkPolicies, PDB, quotas, limites et probes ;
  \item ajout de la supply chain avec SBOM, signatures et promotion ;
  \item ajout des scripts de validation et de support pack ;
  \item alignement de la documentation avec l'etat reel.
\end{enumerate}

Cette creation progressive donne au projet une trajectoire claire : partir d'un
demonstrateur local stable, puis evoluer vers une plateforme production-like
avec controles de securite, haute disponibilite statique, preuves runtime et
release evidence.

\subsection{Legende des diagrammes}

\tikzset{
  srStep/.style={
    rectangle,
    rounded corners=2pt,
    draw=black,
    thick,
    align=center,
    minimum height=8mm,
    text width=31mm,
    fill=gray!10
  },
  srGate/.style={
    diamond,
    draw=black,
    thick,
    align=center,
    aspect=2,
    text width=24mm,
    fill=yellow!15
  },
  srSec/.style={
    rectangle,
    rounded corners=2pt,
    draw=black,
    thick,
    align=center,
    minimum height=8mm,
    text width=32mm,
    fill=red!10
  },
  srProof/.style={
    rectangle,
    rounded corners=2pt,
    draw=black,
    thick,
    align=center,
    minimum height=8mm,
    text width=35mm,
    fill=green!10
  },
  srRuntime/.style={
    rectangle,
    rounded corners=2pt,
    draw=black,
    thick,
    align=center,
    minimum height=8mm,
    text width=34mm,
    fill=blue!8
  },
  srArrow/.style={-{Latex[length=2mm]}, thick}
}

Dans les figures suivantes, les blocs gris representent les etapes techniques,
les blocs rouges les controles de securite, les losanges les gates bloquants,
les blocs bleus les elements runtime Kubernetes et les blocs verts les preuves
archivees.

\subsection{Vue globale de la chaine DevSecOps}

La chaine globale suit une logique progressive : le code est valide en CI, les
images sont construites et securisees, la release est promue par digest, le
deploiement Kubernetes est applique sans reconstruction, puis les preuves
runtime et release sont collectees.

\begin{figure}[htbp]
\centering
\resizebox{\textwidth}{!}{%
\begin{tikzpicture}[node distance=12mm and 10mm]
  \node[srStep] (git) {Depot Git\\monorepo};
  \node[srStep, right=of git] (ci) {Jenkins CI\\qualite et securite};
  \node[srSec, right=of ci] (scan) {SAST, secrets,\\dependances};
  \node[srStep, right=of scan] (build) {Build images\\OCI};
  \node[srSec, right=of build] (supply) {SBOM, signature,\\verification};
  \node[srGate, right=of supply] (gate) {Release\\gate};
  \node[srRuntime, right=of gate] (k8s) {Kubernetes\\production};
  \node[srProof, right=of k8s] (proof) {Preuves\\runtime et release};

  \draw[srArrow] (git) -- (ci);
  \draw[srArrow] (ci) -- (scan);
  \draw[srArrow] (scan) -- (build);
  \draw[srArrow] (build) -- (supply);
  \draw[srArrow] (supply) -- (gate);
  \draw[srArrow] (gate) -- (k8s);
  \draw[srArrow] (k8s) -- (proof);
\end{tikzpicture}}
\caption{Vue globale de la chaine DevSecOps de SecureRAG Hub}
\label{fig:devsecops-global}
\end{figure}

\subsection{Pas 1 : gouvernance du depot et source de verite}

Le premier pas DevSecOps consiste a rendre le depot lisible et gouvernable. Les
sources applicatives, les scripts, l'infrastructure, les runbooks et les
preuves sont separes dans des dossiers dedies. Jenkins est conserve comme
autorite CI/CD officielle, tandis que GitHub Actions reste un heritage
technique non prioritaire.

\begin{figure}[htbp]
\centering
\resizebox{0.95\textwidth}{!}{%
\begin{tikzpicture}[node distance=10mm and 12mm]
  \node[srStep] (repo) {Monorepo\\SecureRAG Hub};
  \node[srStep, below left=of repo] (app) {platform/\\services-laravel/};
  \node[srStep, below=of repo] (infra) {infra/\\k8s, kind, jenkins};
  \node[srStep, below right=of repo] (scripts) {scripts/\\ci, release, deploy};
  \node[srStep, right=of scripts] (docs) {docs/\\runbooks, memoire};
  \node[srProof, right=of repo] (truth) {Source de verite\\Jenkins + docs};

  \draw[srArrow] (repo) -- (app);
  \draw[srArrow] (repo) -- (infra);
  \draw[srArrow] (repo) -- (scripts);
  \draw[srArrow] (repo) -- (docs);
  \draw[srArrow] (repo) -- (truth);
\end{tikzpicture}}
\caption{Organisation du depot et source de verite DevSecOps}
\label{fig:devsecops-repo}
\end{figure}

\subsection{Pas 2 : integration continue Jenkins}

La CI verifie que le code peut etre livre sans introduire de regression
evidente. Elle execute les tests, les validations statiques, les audits de
dependances et les controles de securite. Les resultats sont archives sous
forme de rapports exploitables.

\begin{figure}[htbp]
\centering
\resizebox{\textwidth}{!}{%
\begin{tikzpicture}[node distance=11mm and 9mm]
  \node[srStep] (push) {Push ou\\merge request};
  \node[srStep, right=of push] (checkout) {Checkout\\Jenkins};
  \node[srStep, right=of checkout] (lint) {Lint et rendu\\Kustomize};
  \node[srStep, right=of lint] (tests) {Tests Laravel\\et coverage};
  \node[srSec, right=of tests] (security) {Semgrep\\Gitleaks\\Trivy FS};
  \node[srGate, right=of security] (gate) {Quality\\Gate CI};
  \node[srProof, right=of gate] (reports) {Rapports CI\\archives};

  \draw[srArrow] (push) -- (checkout);
  \draw[srArrow] (checkout) -- (lint);
  \draw[srArrow] (lint) -- (tests);
  \draw[srArrow] (tests) -- (security);
  \draw[srArrow] (security) -- (gate);
  \draw[srArrow] (gate) -- (reports);
\end{tikzpicture}}
\caption{Pipeline Jenkins CI}
\label{fig:devsecops-ci}
\end{figure}

\subsection{Pas 3 : controles de securite applicative et dependances}

Les controles applicatifs sont positionnes tot dans la chaine. Ils incluent la
recherche de secrets, les scans de dependances Composer et npm lorsque les
lockfiles sont presents, les regles SAST et les validations Sonar lorsque
l'environnement fournit les variables necessaires.

\begin{figure}[htbp]
\centering
\resizebox{0.98\textwidth}{!}{%
\begin{tikzpicture}[node distance=10mm and 10mm]
  \node[srStep] (code) {Code Laravel\\et scripts};
  \node[srSec, below left=of code] (deps) {Composer audit\\npm audit};
  \node[srSec, below=of code] (sast) {Semgrep\\Sonar optionnel};
  \node[srSec, below right=of code] (secrets) {Gitleaks\\secrets};
  \node[srGate, right=of code] (block) {Gate\\bloquant};
  \node[srProof, right=of block] (evidence) {security/reports\\artifacts/ci};

  \draw[srArrow] (code) -- (deps);
  \draw[srArrow] (code) -- (sast);
  \draw[srArrow] (code) -- (secrets);
  \draw[srArrow] (deps) -- (block);
  \draw[srArrow] (sast) -- (block);
  \draw[srArrow] (secrets) -- (block);
  \draw[srArrow] (block) -- (evidence);
\end{tikzpicture}}
\caption{Controles de securite applicative et dependances}
\label{fig:devsecops-security-checks}
\end{figure}

\subsection{Pas 4 : construction des images OCI}

Le build produit les images officielles des workloads Laravel. Le runtime
officiel est separe des dossiers legacy afin de ne pas construire ni deployer
des composants sans source applicative exploitable. Les images sont poussees
dans un registre local pour etre consommees par le cluster kind.

\begin{figure}[htbp]
\centering
\resizebox{\textwidth}{!}{%
\begin{tikzpicture}[node distance=11mm and 9mm]
  \node[srStep] (src) {Sources Laravel\\officielles};
  \node[srStep, right=of src] (dockerfile) {Dockerfiles\\production};
  \node[srStep, right=of dockerfile] (build) {docker build};
  \node[srStep, right=of build] (tag) {Tags\\dev/demo/production};
  \node[srStep, right=of tag] (registry) {Registre local\\localhost:5001};
  \node[srRuntime, right=of registry] (kind) {Cluster kind\\pull images};

  \draw[srArrow] (src) -- (dockerfile);
  \draw[srArrow] (dockerfile) -- (build);
  \draw[srArrow] (build) -- (tag);
  \draw[srArrow] (tag) -- (registry);
  \draw[srArrow] (registry) -- (kind);
\end{tikzpicture}}
\caption{Construction et publication des images OCI}
\label{fig:devsecops-build-images}
\end{figure}

\subsection{Pas 5 : supply chain securisee}

La supply chain vise a prouver qu'une image livree correspond bien a l'artefact
construit, scanne, documente par SBOM, signe et verifie. Le projet suit une
logique \textit{digest-first} pour eviter de reconstruire une image au moment du
deploiement.

\begin{figure}[htbp]
\centering
\resizebox{\textwidth}{!}{%
\begin{tikzpicture}[node distance=10mm and 8mm]
  \node[srStep] (image) {Image candidate};
  \node[srSec, right=of image] (trivy) {Trivy\\image scan};
  \node[srSec, right=of trivy] (sbom) {Syft\\SBOM};
  \node[srSec, right=of sbom] (sign) {Cosign\\sign};
  \node[srSec, right=of sign] (verify) {Cosign\\verify};
  \node[srGate, right=of verify] (gate) {Evidence\\gate};
  \node[srStep, right=of gate] (promote) {Promotion\\par digest};
  \node[srProof, right=of promote] (release) {Attestation\\release};

  \draw[srArrow] (image) -- (trivy);
  \draw[srArrow] (trivy) -- (sbom);
  \draw[srArrow] (sbom) -- (sign);
  \draw[srArrow] (sign) -- (verify);
  \draw[srArrow] (verify) -- (gate);
  \draw[srArrow] (gate) -- (promote);
  \draw[srArrow] (promote) -- (release);
\end{tikzpicture}}
\caption{Supply chain securisee : scan, SBOM, signature, verification et promotion}
\label{fig:devsecops-supply-chain}
\end{figure}

\subsection{Pas 6 : deploiement CD sans reconstruction}

Le CD ne doit pas reconstruire les images. Il verifie les artefacts promus,
applique l'overlay Kustomize choisi et collecte les preuves post-deploiement.
Cette separation limite le risque de divergence entre l'image testee et l'image
deployee.

\begin{figure}[htbp]
\centering
\resizebox{\textwidth}{!}{%
\begin{tikzpicture}[node distance=10mm and 9mm]
  \node[srProof] (digest) {Digest promu\\promotion-digests.txt};
  \node[srGate, right=of digest] (verify) {Verification\\pre-deploy};
  \node[srStep, right=of verify] (kustomize) {Kustomize\\overlay};
  \node[srRuntime, right=of kustomize] (apply) {kubectl apply\\cluster};
  \node[srStep, right=of apply] (smoke) {Smoke tests\\security smoke};
  \node[srProof, right=of smoke] (runtime) {Runtime\\evidence};

  \draw[srArrow] (digest) -- (verify);
  \draw[srArrow] (verify) -- (kustomize);
  \draw[srArrow] (kustomize) -- (apply);
  \draw[srArrow] (apply) -- (smoke);
  \draw[srArrow] (smoke) -- (runtime);
\end{tikzpicture}}
\caption{Deploiement CD sans reconstruction d'image}
\label{fig:devsecops-cd-no-rebuild}
\end{figure}

\subsection{Pas 7 : architecture Kubernetes production-like}

L'overlay \texttt{production} conserve le perimetre officiel Laravel et ajoute
les controles attendus pour une plateforme production-like : plusieurs replicas,
rolling update controle, PDB, anti-affinity, topology spread constraints, HPA,
NetworkPolicies et exposition officielle du portail par un Service Kubernetes.

\begin{figure}[htbp]
\centering
\resizebox{\textwidth}{!}{%
\begin{tikzpicture}[node distance=10mm and 10mm]
  \node[srRuntime] (ns) {Namespace\\securerag-hub};
  \node[srRuntime, below left=of ns] (portal) {portal-web\\replicas 3};
  \node[srRuntime, below=of ns] (services) {Services Laravel\\replicas 2};
  \node[srRuntime, below right=of ns] (network) {NetworkPolicies\\default deny};
  \node[srSec, right=of ns] (hardening) {PSA restricted\\securityContext};
  \node[srRuntime, right=of hardening] (ha) {PDB, HPA\\spread};
  \node[srProof, right=of ha] (proof) {production\\runtime evidence};

  \draw[srArrow] (ns) -- (portal);
  \draw[srArrow] (ns) -- (services);
  \draw[srArrow] (ns) -- (network);
  \draw[srArrow] (ns) -- (hardening);
  \draw[srArrow] (hardening) -- (ha);
  \draw[srArrow] (ha) -- (proof);
\end{tikzpicture}}
\caption{Overlay Kubernetes production-like}
\label{fig:devsecops-k8s-production}
\end{figure}

\subsection{Pas 8 : metrics-server et HPA}

Le HPA n'est credible que si \texttt{metrics-server} est disponible et si les
valeurs CPU/memoire ne restent pas a \texttt{\string<unknown\string>}. Le projet distingue
donc les HPA rendus statiquement et les HPA prouves runtime.

\begin{figure}[htbp]
\centering
\resizebox{0.98\textwidth}{!}{%
\begin{tikzpicture}[node distance=10mm and 10mm]
  \node[srRuntime] (cluster) {Cluster\\actif};
  \node[srRuntime, right=of cluster] (metrics) {metrics-server\\Available};
  \node[srStep, right=of metrics] (top) {kubectl top\\nodes/pods};
  \node[srRuntime, right=of top] (hpa) {HPA\\targets remplis};
  \node[srGate, right=of hpa] (gate) {Aucun\\\texttt{\string<unknown\string>}};
  \node[srProof, right=of gate] (report) {hpa-runtime\\report};

  \draw[srArrow] (cluster) -- (metrics);
  \draw[srArrow] (metrics) -- (top);
  \draw[srArrow] (top) -- (hpa);
  \draw[srArrow] (hpa) -- (gate);
  \draw[srArrow] (gate) -- (report);
\end{tikzpicture}}
\caption{Validation runtime de metrics-server et des HPA}
\label{fig:devsecops-hpa}
\end{figure}

\subsection{Pas 9 : Kyverno Audit puis preparation Enforce}

Kyverno est utilise pour auditer les exigences de securite Kubernetes :
profil pod restreint, absence de privileges, references d'images, exposition
des Services et verification Cosign. Le mode \texttt{Enforce} n'est prepare
qu'apres preuve de la supply chain.

\begin{figure}[htbp]
\centering
\resizebox{\textwidth}{!}{%
\begin{tikzpicture}[node distance=10mm and 9mm]
  \node[srRuntime] (cluster) {Cluster\\Kubernetes};
  \node[srSec, right=of cluster] (kyverno) {Kyverno\\Audit};
  \node[srStep, right=of kyverno] (policies) {ClusterPolicies\\SecureRAG};
  \node[srProof, right=of policies] (reports) {PolicyReports};
  \node[srGate, right=of reports] (ready) {Supply chain\\prouvee ?};
  \node[srSec, right=of ready] (enforce) {Enforce\\prudent};

  \draw[srArrow] (cluster) -- (kyverno);
  \draw[srArrow] (kyverno) -- (policies);
  \draw[srArrow] (policies) -- (reports);
  \draw[srArrow] (reports) -- (ready);
  \draw[srArrow] (ready) -- node[above]{oui} (enforce);
\end{tikzpicture}}
\caption{Passage progressif de Kyverno Audit vers Enforce}
\label{fig:devsecops-kyverno}
\end{figure}

\subsection{Pas 10 : gestion des secrets}

La gestion des secrets repose sur une separation entre les valeurs exemples,
les secrets locaux generes, les credentials Jenkins et les secrets Kubernetes.
Aucun secret reel ne doit etre versionne. La rotation reste documentee et doit
etre prouvee dans un environnement d'exploitation.

\begin{figure}[htbp]
\centering
\resizebox{0.98\textwidth}{!}{%
\begin{tikzpicture}[node distance=10mm and 10mm]
  \node[srSec] (strategy) {Strategie\\secrets};
  \node[srStep, below left=of strategy] (examples) {.env.example\\sans secret reel};
  \node[srStep, below=of strategy] (local) {Bootstrap\\local};
  \node[srStep, below right=of strategy] (jenkins) {Credentials\\Jenkins};
  \node[srRuntime, right=of local] (k8s) {Secrets\\Kubernetes};
  \node[srProof, right=of k8s] (audit) {Runbook\\rotation};

  \draw[srArrow] (strategy) -- (examples);
  \draw[srArrow] (strategy) -- (local);
  \draw[srArrow] (strategy) -- (jenkins);
  \draw[srArrow] (local) -- (k8s);
  \draw[srArrow] (jenkins) -- (k8s);
  \draw[srArrow] (k8s) -- (audit);
\end{tikzpicture}}
\caption{Gestion des secrets DevSecOps}
\label{fig:devsecops-secrets}
\end{figure}

\subsection{Pas 11 : resilience donnees, backup et restore}

La resilience donnees production-like impose de sortir d'une logique de
stockage local fragile. L'approche cible est d'utiliser une base externe ou
dediee, de produire des backups, de tester le restore et d'archiver la preuve.

\begin{figure}[htbp]
\centering
\resizebox{\textwidth}{!}{%
\begin{tikzpicture}[node distance=10mm and 9mm]
  \node[srRuntime] (app) {Workloads\\Laravel};
  \node[srRuntime, right=of app] (db) {DB externe\\PostgreSQL};
  \node[srStep, right=of db] (backup) {Backup\\planifie};
  \node[srGate, right=of backup] (restore) {Restore\\teste};
  \node[srProof, right=of restore] (evidence) {Preuve\\data resilience};

  \draw[srArrow] (app) -- (db);
  \draw[srArrow] (db) -- (backup);
  \draw[srArrow] (backup) -- (restore);
  \draw[srArrow] (restore) -- (evidence);
\end{tikzpicture}}
\caption{Resilience donnees : DB externe, backup et restore}
\label{fig:devsecops-data-resilience}
\end{figure}

\subsection{Pas 12 : observabilite runtime}

L'observabilite minimale exploitable repose sur les objets Kubernetes, les
events, les logs applicatifs, les HPA, les metriques \texttt{kubectl top}, les
rapports Kyverno et les snapshots collectes dans les artefacts. Une stack
Prometheus/Grafana/Loki reste une evolution possible, mais n'est pas requise
pour ne pas surcharger la demonstration.

\begin{figure}[htbp]
\centering
\resizebox{\textwidth}{!}{%
\begin{tikzpicture}[node distance=10mm and 8mm]
  \node[srRuntime] (cluster) {Cluster\\runtime};
  \node[srStep, right=of cluster] (objects) {deploy, pods\\svc, pdb, hpa};
  \node[srStep, right=of objects] (events) {Events\\logs};
  \node[srStep, right=of events] (metrics) {kubectl top\\metrics};
  \node[srStep, right=of metrics] (kyverno) {Kyverno\\reports};
  \node[srProof, right=of kyverno] (snapshot) {Observability\\snapshot};

  \draw[srArrow] (cluster) -- (objects);
  \draw[srArrow] (objects) -- (events);
  \draw[srArrow] (events) -- (metrics);
  \draw[srArrow] (metrics) -- (kyverno);
  \draw[srArrow] (kyverno) -- (snapshot);
\end{tikzpicture}}
\caption{Observabilite runtime exploitable}
\label{fig:devsecops-observability}
\end{figure}

\subsection{Pas 13 : support pack et preuves finales}

La derniere etape consiste a consolider les preuves produites par la CI, la
release, Kubernetes, Kyverno, metrics-server, les HPA et les runbooks dans un
support pack final. Ce support pack evite de presenter oralement des controles
non prouves.

\begin{figure}[htbp]
\centering
\resizebox{\textwidth}{!}{%
\begin{tikzpicture}[node distance=10mm and 8mm]
  \node[srProof] (ci) {Preuves CI};
  \node[srProof, right=of ci] (release) {Preuves\\release};
  \node[srProof, right=of release] (runtime) {Preuves\\runtime};
  \node[srProof, right=of runtime] (security) {Preuves\\securite};
  \node[srStep, right=of security] (pack) {Support\\pack};
  \node[srGate, right=of pack] (truth) {Tableau\\etat final};

  \draw[srArrow] (ci) -- (release);
  \draw[srArrow] (release) -- (runtime);
  \draw[srArrow] (runtime) -- (security);
  \draw[srArrow] (security) -- (pack);
  \draw[srArrow] (pack) -- (truth);
\end{tikzpicture}}
\caption{Consolidation des preuves dans le support pack}
\label{fig:devsecops-support-pack}
\end{figure}

\subsection{Synthese des pas DevSecOps}

\begin{center}
\renewcommand{\arraystretch}{1.2}
\begin{tabular}{|p{8mm}|p{42mm}|p{55mm}|p{35mm}|}
\hline
\textbf{Pas} & \textbf{Objectif} & \textbf{Preuve attendue} & \textbf{Etat cible} \\
\hline
1 & Gouvernance du depot & Structure repo, Jenkins source de verite & Termine \\
\hline
2 & CI Jenkins & Rapports tests, lint, scans & Termine si artefacts presents \\
\hline
3 & Securite code et dependances & Semgrep, Gitleaks, Composer/npm audit, Sonar si configure & Termine ou dependant environnement \\
\hline
4 & Build images & Images OCI officielles dans registre local & Termine si Docker/registry actifs \\
\hline
5 & Supply chain & Trivy, Syft, Cosign sign/verify, digest & Termine seulement avec preuves release \\
\hline
6 & CD sans rebuild & Deploiement d'image promue par digest & Termine avec preuve no-rebuild \\
\hline
7 & Kubernetes production-like & Overlay production, PDB, HPA, NetworkPolicies & Termine statique, runtime a prouver \\
\hline
8 & metrics-server et HPA & kubectl top, HPA sans \texttt{\string<unknown\string>} & Dependant cluster actif \\
\hline
9 & Kyverno Audit/Enforce & CRDs, policies, PolicyReports & Audit d'abord, Enforce prudent \\
\hline
10 & Secrets & Aucun secret reel dans Git, bootstrap controle & Partiel sans rotation prouvee \\
\hline
11 & Data resilience & Backup, restore, preuve minimale & Pret non execute si DB absente \\
\hline
12 & Observabilite & Snapshot runtime, logs, events, metrics & Dependant cluster actif \\
\hline
13 & Support pack & Archive finale des preuves & Termine si artefacts generes \\
\hline
\end{tabular}
\end{center}

\subsection{Formulation recommandee pour le memoire}

La partie DevSecOps de \textit{SecureRAG Hub} ne se limite pas a une execution
CI/CD. Elle constitue une chaine de confiance complete reliant le depot Git, la
CI Jenkins, les controles de securite, la construction d'images OCI, la supply
chain par SBOM et signature, le deploiement Kubernetes production-like, les
policies Kyverno, l'observabilite runtime et l'archivage des preuves. Les
controles ne sont declares termines que lorsque les artefacts correspondants
sont presents. Les controles dependant de Docker, kind, metrics-server, Kyverno,
Cosign, Syft, Trivy, PostgreSQL ou Jenkins sont explicitement distingues afin
d'eviter toute surdeclaration.

%
% Packages requis dans le preambule du document principal :
% \usepackage{tikz}
% \usepackage{graphicx}
% \usepackage{array}
% \usetikzlibrary{arrows.meta,positioning,shapes.geometric,fit,backgrounds}

\section{Description complete de la partie DevSecOps}
\label{sec:devsecops-diagrammes-pas-a-pas}

Cette section presente la partie DevSecOps de \textit{SecureRAG Hub} sous la
forme d'un enchainement de pas techniques. L'objectif est de montrer comment le
projet passe d'un code applicatif Laravel versionne a une plateforme
production-like deployee sur Kubernetes, securisee, auditee et documentee par
des preuves.

Le perimetre de cette description est volontairement limite a la chaine
DevSecOps, a Kubernetes, a la supply chain, aux controles de securite, a
l'observabilite, aux preuves runtime et a la documentation. Les composants
metiers hors DevSecOps ne sont pas detailles ici.

\subsection{Creation progressive du socle DevSecOps}
\label{subsec:creation-socle-devsecops}

La creation de la partie DevSecOps n'a pas ete abordee comme un simple ajout
d'outils autour du projet. Elle a ete construite comme un systeme complet de
livraison, de securisation et de preuve. La demarche a consiste a partir du
besoin de soutenance et d'exploitation locale, puis a transformer ce besoin en
une chaine technique reproductible.

Le principe directeur a ete le suivant : chaque automatisation ajoutee au
projet devait produire soit un artefact deployable, soit un rapport de
validation, soit une preuve exploitable dans le memoire. Cette regle a permis
d'eviter une accumulation d'outils non relies entre eux.

\subsubsection{Objectifs initiaux de creation}

Les objectifs initiaux de la creation DevSecOps etaient les suivants :

\begin{itemize}
  \item disposer d'un depot organise et lisible ;
  \item automatiser les tests et les controles de securite ;
  \item construire des images OCI de maniere reproductible ;
  \item deployer localement sur Kubernetes sans infrastructure cloud ;
  \item separer les environnements \texttt{dev}, \texttt{demo} et
  \texttt{production} ;
  \item produire des preuves techniques reutilisables dans la soutenance ;
  \item distinguer honnetement les controles termines des controles dependants
  de l'environnement.
\end{itemize}

\subsubsection{Construction par couches}

La creation s'est faite par couches successives. Chaque couche ajoute un niveau
de maturite supplementaire sans remettre en cause les couches precedentes.

\begin{figure}[htbp]
\centering
\resizebox{\textwidth}{!}{%
\begin{tikzpicture}[node distance=9mm and 8mm]
  \node[srStep] (repo) {Couche 1\\Depot et structure};
  \node[srStep, right=of repo] (ci) {Couche 2\\CI et tests};
  \node[srSec, right=of ci] (security) {Couche 3\\Scans securite};
  \node[srStep, right=of security] (images) {Couche 4\\Images OCI};
  \node[srRuntime, right=of images] (k8s) {Couche 5\\Kubernetes};
  \node[srSec, right=of k8s] (supply) {Couche 6\\Supply chain};
  \node[srProof, right=of supply] (proofs) {Couche 7\\Preuves};

  \draw[srArrow] (repo) -- (ci);
  \draw[srArrow] (ci) -- (security);
  \draw[srArrow] (security) -- (images);
  \draw[srArrow] (images) -- (k8s);
  \draw[srArrow] (k8s) -- (supply);
  \draw[srArrow] (supply) -- (proofs);
\end{tikzpicture}}
\caption{Creation progressive du socle DevSecOps par couches}
\label{fig:creation-devsecops-couches}
\end{figure}

La premiere couche a consiste a stabiliser l'organisation du depot. Les
dossiers \texttt{scripts/}, \texttt{infra/}, \texttt{docs/}, \texttt{security/},
\texttt{platform/} et \texttt{services-laravel/} donnent une responsabilite
claire a chaque partie. Cette separation est importante car elle rend possible
l'automatisation sans melanger code applicatif, infrastructure, runbooks et
preuves.

La deuxieme couche a introduit les scripts de CI. L'objectif etait d'obtenir un
premier signal de qualite avant de construire les images. Les controles de base
sont regroupes dans \texttt{scripts/ci/} et exposes dans le \texttt{Makefile}.
Cette couche verifie notamment les tests, les audits de dependances, les
regles Sonar et les validations statiques.

La troisieme couche a ajoute les controles de securite. Les outils comme
\texttt{Semgrep}, \texttt{Gitleaks} et \texttt{Trivy} permettent de detecter
des erreurs avant la phase de deploiement. Le choix methodologique est de
placer la securite le plus tot possible dans la chaine, afin d'eviter que les
faiblesses ne soient decouvertes uniquement en runtime.

La quatrieme couche a porte sur la construction des images. Les Dockerfiles et
les scripts de build ont ete alignes sur le runtime Laravel officiel. Les
anciens composants legacy sans sources applicatives exploitables ne sont pas
presentes comme build officiel. Ce point est essentiel pour garder une chaine
de preuve honnete.

La cinquieme couche a introduit Kubernetes avec \texttt{kind} et Kustomize. Le
choix de \texttt{kind} permet de simuler localement un environnement
Kubernetes, tandis que Kustomize permet de separer la base commune des overlays
\texttt{dev}, \texttt{demo} et \texttt{production}.

La sixieme couche a ajoute la supply chain : generation de SBOM, scan d'image,
signature Cosign, verification, promotion par digest et deploiement sans
reconstruction. Cette couche transforme le projet d'un simple deploiement local
vers une logique de confiance sur les artefacts.

La septieme couche a formalise la preuve. Les scripts de validation et les
support packs produisent des fichiers sous \texttt{artifacts/}. Ces fichiers
servent a justifier l'etat reel du projet et evitent de declarer un controle
comme termine sans execution correspondante.

\subsubsection{Decisions de conception DevSecOps}

Plusieurs decisions structurantes ont ete prises pendant la creation :

\begin{itemize}
  \item \textbf{Jenkins comme autorite CI/CD} : Jenkins est retenu comme chaine
  officielle pour eviter une double source de verite avec GitHub Actions.
  \item \textbf{Kustomize plutot que manifests dupliques} : la base Kubernetes
  est commune et les differences sont portees par des overlays.
  \item \textbf{Registre local} : \texttt{localhost:5001} permet de demontrer
  la chaine complete sans dependance cloud.
  \item \textbf{Promotion sans rebuild} : une image validee ne doit pas etre
  reconstruite au moment du deploiement.
  \item \textbf{Preuves avant affirmation} : un controle est declare termine
  uniquement si un rapport ou une sortie runtime le prouve.
  \item \textbf{Separation legacy/officiel} : le runtime officiel ne contient
  que les composants Laravel reellement buildables et deployables.
\end{itemize}

\subsubsection{Artefacts crees pendant la construction}

\begin{center}
\renewcommand{\arraystretch}{1.2}
\begin{tabular}{|p{38mm}|p{65mm}|p{42mm}|}
\hline
\textbf{Famille} & \textbf{Role dans la creation DevSecOps} & \textbf{Exemples} \\
\hline
Scripts CI & Executer tests, audits et validations statiques & \texttt{scripts/ci/} \\
\hline
Scripts release & Produire SBOM, signatures, verification et promotion & \texttt{scripts/release/} \\
\hline
Scripts deploy & Creer le cluster, construire les images et deployer & \texttt{scripts/deploy/} \\
\hline
Scripts validate & Collecter les preuves runtime et les rapports & \texttt{scripts/validate/} \\
\hline
Infrastructure & Decrire Jenkins, kind et Kubernetes & \texttt{infra/} \\
\hline
Documentation & Expliquer l'exploitation et les limites reelles & \texttt{docs/runbooks/} \\
\hline
Preuves & Archiver les resultats et statuts reels & \texttt{artifacts/} \\
\hline
\end{tabular}
\end{center}

\subsubsection{Enchainement de creation operationnelle}

L'enchainement operationnel de creation peut etre resume ainsi :

\begin{enumerate}
  \item creation de l'arborescence technique du depot ;
  \item ajout du \texttt{Makefile} comme interface de commandes ;
  \item ajout des scripts CI pour tester et verifier ;
  \item ajout des Dockerfiles et scripts de build ;
  \item ajout du cluster local \texttt{kind} et du registre local ;
  \item creation des manifests Kubernetes de base ;
  \item creation des overlays \texttt{dev}, \texttt{demo} et
  \texttt{production} ;
  \item ajout des NetworkPolicies, PDB, quotas, limites et probes ;
  \item ajout de la supply chain avec SBOM, signatures et promotion ;
  \item ajout des scripts de validation et de support pack ;
  \item alignement de la documentation avec l'etat reel.
\end{enumerate}

Cette creation progressive donne au projet une trajectoire claire : partir d'un
demonstrateur local stable, puis evoluer vers une plateforme production-like
avec controles de securite, haute disponibilite statique, preuves runtime et
release evidence.

\subsection{Legende des diagrammes}

\tikzset{
  srStep/.style={
    rectangle,
    rounded corners=2pt,
    draw=black,
    thick,
    align=center,
    minimum height=8mm,
    text width=31mm,
    fill=gray!10
  },
  srGate/.style={
    diamond,
    draw=black,
    thick,
    align=center,
    aspect=2,
    text width=24mm,
    fill=yellow!15
  },
  srSec/.style={
    rectangle,
    rounded corners=2pt,
    draw=black,
    thick,
    align=center,
    minimum height=8mm,
    text width=32mm,
    fill=red!10
  },
  srProof/.style={
    rectangle,
    rounded corners=2pt,
    draw=black,
    thick,
    align=center,
    minimum height=8mm,
    text width=35mm,
    fill=green!10
  },
  srRuntime/.style={
    rectangle,
    rounded corners=2pt,
    draw=black,
    thick,
    align=center,
    minimum height=8mm,
    text width=34mm,
    fill=blue!8
  },
  srArrow/.style={-{Latex[length=2mm]}, thick}
}

Dans les figures suivantes, les blocs gris representent les etapes techniques,
les blocs rouges les controles de securite, les losanges les gates bloquants,
les blocs bleus les elements runtime Kubernetes et les blocs verts les preuves
archivees.

\subsection{Vue globale de la chaine DevSecOps}

La chaine globale suit une logique progressive : le code est valide en CI, les
images sont construites et securisees, la release est promue par digest, le
deploiement Kubernetes est applique sans reconstruction, puis les preuves
runtime et release sont collectees.

\begin{figure}[htbp]
\centering
\resizebox{\textwidth}{!}{%
\begin{tikzpicture}[node distance=12mm and 10mm]
  \node[srStep] (git) {Depot Git\\monorepo};
  \node[srStep, right=of git] (ci) {Jenkins CI\\qualite et securite};
  \node[srSec, right=of ci] (scan) {SAST, secrets,\\dependances};
  \node[srStep, right=of scan] (build) {Build images\\OCI};
  \node[srSec, right=of build] (supply) {SBOM, signature,\\verification};
  \node[srGate, right=of supply] (gate) {Release\\gate};
  \node[srRuntime, right=of gate] (k8s) {Kubernetes\\production};
  \node[srProof, right=of k8s] (proof) {Preuves\\runtime et release};

  \draw[srArrow] (git) -- (ci);
  \draw[srArrow] (ci) -- (scan);
  \draw[srArrow] (scan) -- (build);
  \draw[srArrow] (build) -- (supply);
  \draw[srArrow] (supply) -- (gate);
  \draw[srArrow] (gate) -- (k8s);
  \draw[srArrow] (k8s) -- (proof);
\end{tikzpicture}}
\caption{Vue globale de la chaine DevSecOps de SecureRAG Hub}
\label{fig:devsecops-global}
\end{figure}

\subsection{Pas 1 : gouvernance du depot et source de verite}

Le premier pas DevSecOps consiste a rendre le depot lisible et gouvernable. Les
sources applicatives, les scripts, l'infrastructure, les runbooks et les
preuves sont separes dans des dossiers dedies. Jenkins est conserve comme
autorite CI/CD officielle, tandis que GitHub Actions reste un heritage
technique non prioritaire.

\begin{figure}[htbp]
\centering
\resizebox{0.95\textwidth}{!}{%
\begin{tikzpicture}[node distance=10mm and 12mm]
  \node[srStep] (repo) {Monorepo\\SecureRAG Hub};
  \node[srStep, below left=of repo] (app) {platform/\\services-laravel/};
  \node[srStep, below=of repo] (infra) {infra/\\k8s, kind, jenkins};
  \node[srStep, below right=of repo] (scripts) {scripts/\\ci, release, deploy};
  \node[srStep, right=of scripts] (docs) {docs/\\runbooks, memoire};
  \node[srProof, right=of repo] (truth) {Source de verite\\Jenkins + docs};

  \draw[srArrow] (repo) -- (app);
  \draw[srArrow] (repo) -- (infra);
  \draw[srArrow] (repo) -- (scripts);
  \draw[srArrow] (repo) -- (docs);
  \draw[srArrow] (repo) -- (truth);
\end{tikzpicture}}
\caption{Organisation du depot et source de verite DevSecOps}
\label{fig:devsecops-repo}
\end{figure}

\subsection{Pas 2 : integration continue Jenkins}

La CI verifie que le code peut etre livre sans introduire de regression
evidente. Elle execute les tests, les validations statiques, les audits de
dependances et les controles de securite. Les resultats sont archives sous
forme de rapports exploitables.

\begin{figure}[htbp]
\centering
\resizebox{\textwidth}{!}{%
\begin{tikzpicture}[node distance=11mm and 9mm]
  \node[srStep] (push) {Push ou\\merge request};
  \node[srStep, right=of push] (checkout) {Checkout\\Jenkins};
  \node[srStep, right=of checkout] (lint) {Lint et rendu\\Kustomize};
  \node[srStep, right=of lint] (tests) {Tests Laravel\\et coverage};
  \node[srSec, right=of tests] (security) {Semgrep\\Gitleaks\\Trivy FS};
  \node[srGate, right=of security] (gate) {Quality\\Gate CI};
  \node[srProof, right=of gate] (reports) {Rapports CI\\archives};

  \draw[srArrow] (push) -- (checkout);
  \draw[srArrow] (checkout) -- (lint);
  \draw[srArrow] (lint) -- (tests);
  \draw[srArrow] (tests) -- (security);
  \draw[srArrow] (security) -- (gate);
  \draw[srArrow] (gate) -- (reports);
\end{tikzpicture}}
\caption{Pipeline Jenkins CI}
\label{fig:devsecops-ci}
\end{figure}

\subsection{Pas 3 : controles de securite applicative et dependances}

Les controles applicatifs sont positionnes tot dans la chaine. Ils incluent la
recherche de secrets, les scans de dependances Composer et npm lorsque les
lockfiles sont presents, les regles SAST et les validations Sonar lorsque
l'environnement fournit les variables necessaires.

\begin{figure}[htbp]
\centering
\resizebox{0.98\textwidth}{!}{%
\begin{tikzpicture}[node distance=10mm and 10mm]
  \node[srStep] (code) {Code Laravel\\et scripts};
  \node[srSec, below left=of code] (deps) {Composer audit\\npm audit};
  \node[srSec, below=of code] (sast) {Semgrep\\Sonar optionnel};
  \node[srSec, below right=of code] (secrets) {Gitleaks\\secrets};
  \node[srGate, right=of code] (block) {Gate\\bloquant};
  \node[srProof, right=of block] (evidence) {security/reports\\artifacts/ci};

  \draw[srArrow] (code) -- (deps);
  \draw[srArrow] (code) -- (sast);
  \draw[srArrow] (code) -- (secrets);
  \draw[srArrow] (deps) -- (block);
  \draw[srArrow] (sast) -- (block);
  \draw[srArrow] (secrets) -- (block);
  \draw[srArrow] (block) -- (evidence);
\end{tikzpicture}}
\caption{Controles de securite applicative et dependances}
\label{fig:devsecops-security-checks}
\end{figure}

\subsection{Pas 4 : construction des images OCI}

Le build produit les images officielles des workloads Laravel. Le runtime
officiel est separe des dossiers legacy afin de ne pas construire ni deployer
des composants sans source applicative exploitable. Les images sont poussees
dans un registre local pour etre consommees par le cluster kind.

\begin{figure}[htbp]
\centering
\resizebox{\textwidth}{!}{%
\begin{tikzpicture}[node distance=11mm and 9mm]
  \node[srStep] (src) {Sources Laravel\\officielles};
  \node[srStep, right=of src] (dockerfile) {Dockerfiles\\production};
  \node[srStep, right=of dockerfile] (build) {docker build};
  \node[srStep, right=of build] (tag) {Tags\\dev/demo/production};
  \node[srStep, right=of tag] (registry) {Registre local\\localhost:5001};
  \node[srRuntime, right=of registry] (kind) {Cluster kind\\pull images};

  \draw[srArrow] (src) -- (dockerfile);
  \draw[srArrow] (dockerfile) -- (build);
  \draw[srArrow] (build) -- (tag);
  \draw[srArrow] (tag) -- (registry);
  \draw[srArrow] (registry) -- (kind);
\end{tikzpicture}}
\caption{Construction et publication des images OCI}
\label{fig:devsecops-build-images}
\end{figure}

\subsection{Pas 5 : supply chain securisee}

La supply chain vise a prouver qu'une image livree correspond bien a l'artefact
construit, scanne, documente par SBOM, signe et verifie. Le projet suit une
logique \textit{digest-first} pour eviter de reconstruire une image au moment du
deploiement.

\begin{figure}[htbp]
\centering
\resizebox{\textwidth}{!}{%
\begin{tikzpicture}[node distance=10mm and 8mm]
  \node[srStep] (image) {Image candidate};
  \node[srSec, right=of image] (trivy) {Trivy\\image scan};
  \node[srSec, right=of trivy] (sbom) {Syft\\SBOM};
  \node[srSec, right=of sbom] (sign) {Cosign\\sign};
  \node[srSec, right=of sign] (verify) {Cosign\\verify};
  \node[srGate, right=of verify] (gate) {Evidence\\gate};
  \node[srStep, right=of gate] (promote) {Promotion\\par digest};
  \node[srProof, right=of promote] (release) {Attestation\\release};

  \draw[srArrow] (image) -- (trivy);
  \draw[srArrow] (trivy) -- (sbom);
  \draw[srArrow] (sbom) -- (sign);
  \draw[srArrow] (sign) -- (verify);
  \draw[srArrow] (verify) -- (gate);
  \draw[srArrow] (gate) -- (promote);
  \draw[srArrow] (promote) -- (release);
\end{tikzpicture}}
\caption{Supply chain securisee : scan, SBOM, signature, verification et promotion}
\label{fig:devsecops-supply-chain}
\end{figure}

\subsection{Pas 6 : deploiement CD sans reconstruction}

Le CD ne doit pas reconstruire les images. Il verifie les artefacts promus,
applique l'overlay Kustomize choisi et collecte les preuves post-deploiement.
Cette separation limite le risque de divergence entre l'image testee et l'image
deployee.

\begin{figure}[htbp]
\centering
\resizebox{\textwidth}{!}{%
\begin{tikzpicture}[node distance=10mm and 9mm]
  \node[srProof] (digest) {Digest promu\\promotion-digests.txt};
  \node[srGate, right=of digest] (verify) {Verification\\pre-deploy};
  \node[srStep, right=of verify] (kustomize) {Kustomize\\overlay};
  \node[srRuntime, right=of kustomize] (apply) {kubectl apply\\cluster};
  \node[srStep, right=of apply] (smoke) {Smoke tests\\security smoke};
  \node[srProof, right=of smoke] (runtime) {Runtime\\evidence};

  \draw[srArrow] (digest) -- (verify);
  \draw[srArrow] (verify) -- (kustomize);
  \draw[srArrow] (kustomize) -- (apply);
  \draw[srArrow] (apply) -- (smoke);
  \draw[srArrow] (smoke) -- (runtime);
\end{tikzpicture}}
\caption{Deploiement CD sans reconstruction d'image}
\label{fig:devsecops-cd-no-rebuild}
\end{figure}

\subsection{Pas 7 : architecture Kubernetes production-like}

L'overlay \texttt{production} conserve le perimetre officiel Laravel et ajoute
les controles attendus pour une plateforme production-like : plusieurs replicas,
rolling update controle, PDB, anti-affinity, topology spread constraints, HPA,
NetworkPolicies et exposition officielle du portail par un Service Kubernetes.

\begin{figure}[htbp]
\centering
\resizebox{\textwidth}{!}{%
\begin{tikzpicture}[node distance=10mm and 10mm]
  \node[srRuntime] (ns) {Namespace\\securerag-hub};
  \node[srRuntime, below left=of ns] (portal) {portal-web\\replicas 3};
  \node[srRuntime, below=of ns] (services) {Services Laravel\\replicas 2};
  \node[srRuntime, below right=of ns] (network) {NetworkPolicies\\default deny};
  \node[srSec, right=of ns] (hardening) {PSA restricted\\securityContext};
  \node[srRuntime, right=of hardening] (ha) {PDB, HPA\\spread};
  \node[srProof, right=of ha] (proof) {production\\runtime evidence};

  \draw[srArrow] (ns) -- (portal);
  \draw[srArrow] (ns) -- (services);
  \draw[srArrow] (ns) -- (network);
  \draw[srArrow] (ns) -- (hardening);
  \draw[srArrow] (hardening) -- (ha);
  \draw[srArrow] (ha) -- (proof);
\end{tikzpicture}}
\caption{Overlay Kubernetes production-like}
\label{fig:devsecops-k8s-production}
\end{figure}

\subsection{Pas 8 : metrics-server et HPA}

Le HPA n'est credible que si \texttt{metrics-server} est disponible et si les
valeurs CPU/memoire ne restent pas a \texttt{\string<unknown\string>}. Le projet distingue
donc les HPA rendus statiquement et les HPA prouves runtime.

\begin{figure}[htbp]
\centering
\resizebox{0.98\textwidth}{!}{%
\begin{tikzpicture}[node distance=10mm and 10mm]
  \node[srRuntime] (cluster) {Cluster\\actif};
  \node[srRuntime, right=of cluster] (metrics) {metrics-server\\Available};
  \node[srStep, right=of metrics] (top) {kubectl top\\nodes/pods};
  \node[srRuntime, right=of top] (hpa) {HPA\\targets remplis};
  \node[srGate, right=of hpa] (gate) {Aucun\\\texttt{\string<unknown\string>}};
  \node[srProof, right=of gate] (report) {hpa-runtime\\report};

  \draw[srArrow] (cluster) -- (metrics);
  \draw[srArrow] (metrics) -- (top);
  \draw[srArrow] (top) -- (hpa);
  \draw[srArrow] (hpa) -- (gate);
  \draw[srArrow] (gate) -- (report);
\end{tikzpicture}}
\caption{Validation runtime de metrics-server et des HPA}
\label{fig:devsecops-hpa}
\end{figure}

\subsection{Pas 9 : Kyverno Audit puis preparation Enforce}

Kyverno est utilise pour auditer les exigences de securite Kubernetes :
profil pod restreint, absence de privileges, references d'images, exposition
des Services et verification Cosign. Le mode \texttt{Enforce} n'est prepare
qu'apres preuve de la supply chain.

\begin{figure}[htbp]
\centering
\resizebox{\textwidth}{!}{%
\begin{tikzpicture}[node distance=10mm and 9mm]
  \node[srRuntime] (cluster) {Cluster\\Kubernetes};
  \node[srSec, right=of cluster] (kyverno) {Kyverno\\Audit};
  \node[srStep, right=of kyverno] (policies) {ClusterPolicies\\SecureRAG};
  \node[srProof, right=of policies] (reports) {PolicyReports};
  \node[srGate, right=of reports] (ready) {Supply chain\\prouvee ?};
  \node[srSec, right=of ready] (enforce) {Enforce\\prudent};

  \draw[srArrow] (cluster) -- (kyverno);
  \draw[srArrow] (kyverno) -- (policies);
  \draw[srArrow] (policies) -- (reports);
  \draw[srArrow] (reports) -- (ready);
  \draw[srArrow] (ready) -- node[above]{oui} (enforce);
\end{tikzpicture}}
\caption{Passage progressif de Kyverno Audit vers Enforce}
\label{fig:devsecops-kyverno}
\end{figure}

\subsection{Pas 10 : gestion des secrets}

La gestion des secrets repose sur une separation entre les valeurs exemples,
les secrets locaux generes, les credentials Jenkins et les secrets Kubernetes.
Aucun secret reel ne doit etre versionne. La rotation reste documentee et doit
etre prouvee dans un environnement d'exploitation.

\begin{figure}[htbp]
\centering
\resizebox{0.98\textwidth}{!}{%
\begin{tikzpicture}[node distance=10mm and 10mm]
  \node[srSec] (strategy) {Strategie\\secrets};
  \node[srStep, below left=of strategy] (examples) {.env.example\\sans secret reel};
  \node[srStep, below=of strategy] (local) {Bootstrap\\local};
  \node[srStep, below right=of strategy] (jenkins) {Credentials\\Jenkins};
  \node[srRuntime, right=of local] (k8s) {Secrets\\Kubernetes};
  \node[srProof, right=of k8s] (audit) {Runbook\\rotation};

  \draw[srArrow] (strategy) -- (examples);
  \draw[srArrow] (strategy) -- (local);
  \draw[srArrow] (strategy) -- (jenkins);
  \draw[srArrow] (local) -- (k8s);
  \draw[srArrow] (jenkins) -- (k8s);
  \draw[srArrow] (k8s) -- (audit);
\end{tikzpicture}}
\caption{Gestion des secrets DevSecOps}
\label{fig:devsecops-secrets}
\end{figure}

\subsection{Pas 11 : resilience donnees, backup et restore}

La resilience donnees production-like impose de sortir d'une logique de
stockage local fragile. L'approche cible est d'utiliser une base externe ou
dediee, de produire des backups, de tester le restore et d'archiver la preuve.

\begin{figure}[htbp]
\centering
\resizebox{\textwidth}{!}{%
\begin{tikzpicture}[node distance=10mm and 9mm]
  \node[srRuntime] (app) {Workloads\\Laravel};
  \node[srRuntime, right=of app] (db) {DB externe\\PostgreSQL};
  \node[srStep, right=of db] (backup) {Backup\\planifie};
  \node[srGate, right=of backup] (restore) {Restore\\teste};
  \node[srProof, right=of restore] (evidence) {Preuve\\data resilience};

  \draw[srArrow] (app) -- (db);
  \draw[srArrow] (db) -- (backup);
  \draw[srArrow] (backup) -- (restore);
  \draw[srArrow] (restore) -- (evidence);
\end{tikzpicture}}
\caption{Resilience donnees : DB externe, backup et restore}
\label{fig:devsecops-data-resilience}
\end{figure}

\subsection{Pas 12 : observabilite runtime}

L'observabilite minimale exploitable repose sur les objets Kubernetes, les
events, les logs applicatifs, les HPA, les metriques \texttt{kubectl top}, les
rapports Kyverno et les snapshots collectes dans les artefacts. Une stack
Prometheus/Grafana/Loki reste une evolution possible, mais n'est pas requise
pour ne pas surcharger la demonstration.

\begin{figure}[htbp]
\centering
\resizebox{\textwidth}{!}{%
\begin{tikzpicture}[node distance=10mm and 8mm]
  \node[srRuntime] (cluster) {Cluster\\runtime};
  \node[srStep, right=of cluster] (objects) {deploy, pods\\svc, pdb, hpa};
  \node[srStep, right=of objects] (events) {Events\\logs};
  \node[srStep, right=of events] (metrics) {kubectl top\\metrics};
  \node[srStep, right=of metrics] (kyverno) {Kyverno\\reports};
  \node[srProof, right=of kyverno] (snapshot) {Observability\\snapshot};

  \draw[srArrow] (cluster) -- (objects);
  \draw[srArrow] (objects) -- (events);
  \draw[srArrow] (events) -- (metrics);
  \draw[srArrow] (metrics) -- (kyverno);
  \draw[srArrow] (kyverno) -- (snapshot);
\end{tikzpicture}}
\caption{Observabilite runtime exploitable}
\label{fig:devsecops-observability}
\end{figure}

\subsection{Pas 13 : support pack et preuves finales}

La derniere etape consiste a consolider les preuves produites par la CI, la
release, Kubernetes, Kyverno, metrics-server, les HPA et les runbooks dans un
support pack final. Ce support pack evite de presenter oralement des controles
non prouves.

\begin{figure}[htbp]
\centering
\resizebox{\textwidth}{!}{%
\begin{tikzpicture}[node distance=10mm and 8mm]
  \node[srProof] (ci) {Preuves CI};
  \node[srProof, right=of ci] (release) {Preuves\\release};
  \node[srProof, right=of release] (runtime) {Preuves\\runtime};
  \node[srProof, right=of runtime] (security) {Preuves\\securite};
  \node[srStep, right=of security] (pack) {Support\\pack};
  \node[srGate, right=of pack] (truth) {Tableau\\etat final};

  \draw[srArrow] (ci) -- (release);
  \draw[srArrow] (release) -- (runtime);
  \draw[srArrow] (runtime) -- (security);
  \draw[srArrow] (security) -- (pack);
  \draw[srArrow] (pack) -- (truth);
\end{tikzpicture}}
\caption{Consolidation des preuves dans le support pack}
\label{fig:devsecops-support-pack}
\end{figure}

\subsection{Synthese des pas DevSecOps}

\begin{center}
\renewcommand{\arraystretch}{1.2}
\begin{tabular}{|p{8mm}|p{42mm}|p{55mm}|p{35mm}|}
\hline
\textbf{Pas} & \textbf{Objectif} & \textbf{Preuve attendue} & \textbf{Etat cible} \\
\hline
1 & Gouvernance du depot & Structure repo, Jenkins source de verite & Termine \\
\hline
2 & CI Jenkins & Rapports tests, lint, scans & Termine si artefacts presents \\
\hline
3 & Securite code et dependances & Semgrep, Gitleaks, Composer/npm audit, Sonar si configure & Termine ou dependant environnement \\
\hline
4 & Build images & Images OCI officielles dans registre local & Termine si Docker/registry actifs \\
\hline
5 & Supply chain & Trivy, Syft, Cosign sign/verify, digest & Termine seulement avec preuves release \\
\hline
6 & CD sans rebuild & Deploiement d'image promue par digest & Termine avec preuve no-rebuild \\
\hline
7 & Kubernetes production-like & Overlay production, PDB, HPA, NetworkPolicies & Termine statique, runtime a prouver \\
\hline
8 & metrics-server et HPA & kubectl top, HPA sans \texttt{\string<unknown\string>} & Dependant cluster actif \\
\hline
9 & Kyverno Audit/Enforce & CRDs, policies, PolicyReports & Audit d'abord, Enforce prudent \\
\hline
10 & Secrets & Aucun secret reel dans Git, bootstrap controle & Partiel sans rotation prouvee \\
\hline
11 & Data resilience & Backup, restore, preuve minimale & Pret non execute si DB absente \\
\hline
12 & Observabilite & Snapshot runtime, logs, events, metrics & Dependant cluster actif \\
\hline
13 & Support pack & Archive finale des preuves & Termine si artefacts generes \\
\hline
\end{tabular}
\end{center}

\subsection{Formulation recommandee pour le memoire}

La partie DevSecOps de \textit{SecureRAG Hub} ne se limite pas a une execution
CI/CD. Elle constitue une chaine de confiance complete reliant le depot Git, la
CI Jenkins, les controles de securite, la construction d'images OCI, la supply
chain par SBOM et signature, le deploiement Kubernetes production-like, les
policies Kyverno, l'observabilite runtime et l'archivage des preuves. Les
controles ne sont declares termines que lorsque les artefacts correspondants
sont presents. Les controles dependant de Docker, kind, metrics-server, Kyverno,
Cosign, Syft, Trivy, PostgreSQL ou Jenkins sont explicitement distingues afin
d'eviter toute surdeclaration.
